% concatenated from ManuscriptDDI.tex
% \documentclass[12pt]{article}
\documentclass[12pt, a4paper]{amsart}

\usepackage[left=20mm, right=20mm, top=30mm, bottom=30mm, marginparwidth=80pt]{geometry}


\usepackage{bm}
\usepackage[authoryear,longnamesfirst,sort,numbers]{natbib}
\usepackage{amsfonts}
\usepackage{amsthm, amsmath, amssymb}
\usepackage{graphicx}
\usepackage{xcolor}
\usepackage{booktabs} % for the notation table
\usepackage{subcaption} % for figure 4
\usepackage{float}
\usepackage{lipsum}
\usepackage{tikz}
\numberwithin{equation}{section}

\usepackage{enumitem}

%%% Autonumbering equations
\usepackage{mathtools}
\mathtoolsset{showonlyrefs}

%% set margins (set back to 27mm when fisnished reviewing)
\usepackage{hyperref}
\usepackage{bbm}

\usepackage{xifthen}


%%% MACROS
\newcommand{\RR}{\mathbb{R}}
\newcommand{\NN}{\mathbb{N}}
\newcommand{\diag}{\mathrm{diag}}
\newcommand{\ran}{\mathrm{ran}}

\renewcommand{\Re}{\mathrm{Re}\,}

% Constant steady states:
\newcommand{\CSS}{\bar{X}}
\newcommand{\Cun}{\bar{\un}}
\newcommand{\Cus}{\bar{\us}}
\newcommand{\Cuf}{\bar{\uf}}
\newcommand{\Cud}{\bar{\vv}}

% Nonconstant steady states:
\newcommand{\NSS}{\tilde{X}}

\newcommand{\Nun}{\tilde{\un}}
\newcommand{\Nus}{\tilde{\us}}
\newcommand{\Nuf}{\tilde{\uf}}
\newcommand{\Nud}{\tilde{\vv}}

\newcommand{\un}{u}
\newcommand{\us}{v}
\newcommand{\uf}{w}

\newcommand{\ud}{\vv}

\newcommand{\Ds}{D_v}
\newcommand{\Df}{D_w}

\newcommand{\mn}{m_{u}}  %m
\newcommand{\ms}{m_{v}}  %k_s
\newcommand{\mfa}{m_{w}} %k_f
\newcommand{\md}{m_{d}}  %k
\newcommand{\ffull}{F}

\newcommand{\fn}{f}
\newcommand{\fs}{g}
\newcommand{\ffa}{h}

\newcommand{\fd}{\gg}

\def\mf{\mu_1}
\def\mb{\mu_b}
\def\ml{\mu_2}
\def\me{\mu_3}

\def\uu{u}
\def\vv{v}

\def\ff{f}
\let\muchmore\gg
\renewcommand{\gg}{g}

\renewcommand{\L}{{\mathcal{L}}}
\newcommand{\N}{{\mathcal{N}}}  
\def\JODE{\bm{A}_{\ast}}
\def\Jsmall{\hat{J}}

\newcommand{\Li}{{L^\infty(\Omega)}}
\newcommand{\LiR}{{L^\infty(\RR)}}
\newcommand{\W}[2]
{{
		W^{#1,#2}_
		{
			\ifthenelse
			{
				\equal{#1}{2}
			}
			{
				\nu
			}
			{
			}
		}(\Omega)
}}

\newcommand{\Lp}[1]
{{
		L^
		{
			\ifthenelse
			{
				\isempty{#1}
			}
			{
				p
			}
			{
				#1
			}
		}(\Omega)
}}

\definecolor{lime}{HTML}{A6CE39}
\DeclareRobustCommand{\orcidicon}{%
	\begin{tikzpicture}
		\draw[lime, fill=lime] (0,0) 
		circle [radius=0.16] 
		node[white] {{\fontfamily{qag}\selectfont \tiny ID}};
		\draw[white, fill=white] (-0.0625,0.095) 
		circle [radius=0.007];
	\end{tikzpicture}
	\hspace{-2mm}
}

\newenvironment{nalign}
{\begin{equation}\begin{aligned}}
		{\end{aligned}\end{equation}\ignorespacesafterend}


% DEFINITIONS - Maths Environments
\def\be{\begin{equation}}
\def\ee{\end{equation}}
\def\bpm{\begin{pmatrix}}
\def\epm{\end{pmatrix}}


% DEFINITIONS - All types of brackets
\def\lp{\left(}
\def\rp{\right)}

% DEFINITIONS - Usual Operators
\def\del{\partial}
\def\grad{\nabla}
\DeclareMathOperator{\tr}{tr}


%%%%%%%%%%% Defining Enunciations  %%%%%%%%%%%
\newtheorem{theorem}{\bf Theorem}[section]
\newtheorem{condition}[theorem]{\bf Condition}
\newtheorem{corollary}[theorem]{\bf Corollary}
\newtheorem{proposition}[theorem]{\bf Proposition}
\newtheorem{definition}[theorem]{\bf Definition}
\newtheorem{lemma}[theorem]{\bf Lemma}
\newtheorem{remark}[theorem]{\bf Remark}
\newtheorem{example}[theorem]{\bf Example}
\newtheorem{assumption}[theorem]{\bf Assumption}
\newtheorem{hypothesis}[theorem]{\bf Hypothesis}
%%%%%%%%%%%%%%%%%%%%%%%%%%%%%%%%%%%%%%%%%%%%%%%

\title[Multiple diffusion scales and DDI]{Multiple diffusion scales and diffusion-driven instability: Emergence of near- and far-from-equilibrium patterns} 

% \title{Spatial scale separation and emergent patterns in coupled diffusive--non--diffusive systems} 



% \author{Theo André, Szymon Cygan, Anna Marciniak-Czochra, Finn Münnich}\\

\author[T. André]{Théo André \href{https://orcid.org/0009-0001-7320-1941}{\orcidicon}}
\address[T. André]{Institute for Mathematics, Heidelberg University, Im Neuenheimer Feld 205, 69120 Heidelberg, Germany\\
\href{https://orcid.org/0009-0001-7320-1941}{orcid.org/0009-0001-7320-1941}}
\email{theo.andre@uni-heidelberg.de}

\author[S. Cygan]{Szymon Cygan \href{https://orcid.org/0000-0002-8601-829X}{\orcidicon}}
\address[S. Cygan]{Institute for Mathematics, Heidelberg University, Im Neuenheimer Feld 205, 69120 Heidelberg, Germany\\
Instytut Matematyczny, Uniwersytet Wroc\l{}awski, pl. Grunwaldzki 2/4, \hbox{50-384} Wroc\l{}aw, Poland \\ 
\href{https://orcid.org/0000-0002-8601-829X}{orcid.org/0000-0002-8601-829X}}
\email{szymon.cygan@uni-heidelberg.de}
\urladdr {http://www.math.uni.wroc.pl/~scygan}

\author[A. Marciniak-Czochra]{Anna Marciniak-Czochra \href{https://orcid.org/0000-0002-5831-6505}{\orcidicon}}
\address[A. Marciniak-Czochra]{Institute for Mathematics and IWR, Heidelberg University, Im Neuenheimer Feld 205, 69120 Heidelberg, Germany\\ 
\href{https://orcid.org/0000-0002-5831-6505}{orcid.org/0000-0002-5831-6505}}
\email{anna.marciniak@iwr.uni-heidelberg.de}
\urladdr{\href{https://biostruct.iwr.uni-heidelberg.de/folder_people/Anna.Marciniak/index.html}{https://biostruct.iwr.uni-heidelberg.de/folder\_people/Anna.Marciniak/index.html}}

\author[F. Münnich]{Finn Münnich \href{https://orcid.org/0009-0007-9384-2002}{\orcidicon}}
\address[F. Münnich]{Institute for Mathematics, Heidelberg University, Im Neuenheimer Feld 205, 69120 Heidelberg, Germany\\
\href{https://orcid.org/0009-0007-9384-2002}{orcid.org/0009-0007-9384-2002}}
\email{finn.muennich@stud.uni-heidelberg.de}

% Institute for Mathematics, Heidelberg University, Im Neuenheimer Feld 205, 69120 Heidelberg, Germany\\
% Institute for Mathematics and IWR, Heidelberg University, Im Neuenheimer Feld 205, 69120 Heidelberg, Germany

\date{}
%%%%%%%%%%%%%%%%%%%%%%%%%%%%%%%%%%%%%%%%%%%%%%%


\begin{document}

%%%% Keyword entries to be placed here %%%%
\keywords{Reaction--diffusion equation, pattern formation, diffusion-driven instability, discontinuous steady states}



%%%% Abstract text to be placed here %%%%%%%%%%%%
\begin{abstract}
    This paper investigates pattern formation in reaction--diffusion systems with both diffusive and nondiffusive components, providing necessary and sufficient conditions for diffusion-driven instability (DDI) and establishing the existence of far-from-equilibrium patterns.
    While previous work has linked DDI to instability in the purely nondiffusive subsystem -- thereby destabilizing all regular Turing patterns -- we show that DDI can also arise from subsystems involving nondiffusive and slow-diffusive components using three different spatial scales. This leads to simple sufficient conditions for DDI in systems with arbitrary numbers of components.
    Moreover, we fully classify all possible sources of DDI in the case of two diffusive and one nondiffusive component.
    Further, we prove the existence of far-from-equilibrium patterns exhibiting branch-switching and discontinuities in the nondiffusive components, which cannot occur in classical reaction--diffusion equations. 
    We illustrate our results with a receptor-based model supported by numerical bifurcation analysis and simulations.
    These findings extend the theoretical foundations of pattern formation, demonstrating how coupling between diffusive and nondiffusive dynamics can generate patterns beyond the reach of the classical reaction--diffusion framework.
\end{abstract}

\maketitle

% \begin{enumerate}
%     \item {\color{gray} Why RD-ODE? Examples from Ecology, cancer, ligands receptors, Hydra}
%     \item {\color{gray}We are interested to spontaneous pattern formation and de-novo patterns. Linked to DDI}
%     \item {\color{gray}A lot of results for RD-ODE systems with single diffusive component.}
%     \item {\color{gray}Well known to have DDI in RD-ODE systems (with one diffusion) via autocatalysis condition in nondiffusing components}
%     \item {\color{gray}But if only one diffusive component, all Turing patterns unstable}
%     \item {\color{gray}Need at least two diffusive terms to observe regular patterns}
%     \item {\color{gray} From application, the diffusive components often have different diffusion rates. }
%     \item {\color{gray}so assume we can split the rd-ode model in 3 parts: non, slow, fast}
%     \item {\color{gray} How does DDI look like in systems of this structure?}
%     \item {\color{gray}Autocatalysis condition can be transferred but destabilizes all nearby TPs}
%     \item {\color{gray} Aim: Derive simple conditions for DDI (easy verifiable) in a RD-ODE system assuming the mentioned structure, under which TP possible. (Make clear that this is one of main points of this work)}
%     \item {\color{gray}Under these assumptions TP possible (cite?) but if TP unstable, we can still observe jump patterns (far-from-equilibrium) similar to RD-ODE systems with one diffusive components}
%     \item {\color{gray}The existence result of jump patterns can be generalized to multiple diffusive components (and is actually independent of size of diffusion)}
%     \item {\color{gray} We state the generalized proofs for readers convenience}
% \end{enumerate}

% Section 1 : Introduction
\section{Introduction}\label{sec:intro}

% {\color{gray}

% Beyond classical Turing patterns, reaction--diffusion--ODE systems are also known to exhibit far-from-equilibrium spatial structures, including spike-like solutions and jump discontinuities \cite{Haerting17, Marciniak15}. These patterns arise through mechanisms that differ fundamentally from DDI and may persist even in parameter regimes where Turing patterns are unstable.

% }

A central question in pattern formation concerns the emergence of spatial structure from spatially homogeneous steady states. One of the primary mechanisms underlying this phenomenon is diffusion-driven instability (DDI), also known as Turing instability \cite{turing1952chemical}. In classical reaction–diffusion systems, DDI is a bifurcation from a stable homogeneous steady state with respect to diffusion parameters, occurring when the effective diffusion ratio exceeds a critical threshold \cite{Nishiura1982}. Near onset, this instability gives rise to Turing patterns, typically small-amplitude, spatially heterogeneous solutions near the homogeneous equilibrium. In contrast, coupling even a single reaction–diffusion equation with non-diffusive dynamics, described by space-dependent ODEs, leads to fundamentally different behaviour. In such systems,  DDI may arise through an autocatalytic mechanism in the non-diffusive subsystem, which both ensures the onset of instability of the spatially constant solution and implies instability of all Turing patterns \cite{Marciniak15,Cygan2022}. As a consequence, observable patterns are necessarily far-from-equilibrium and exhibit singular structures such as spikes \cite{HMC14} or jump discontinuities \cite{Haerting17}.

This raises the question of whether generic reaction-diffusion-ODE systems with multiple diffusion scales also exhibit such degenerate dynamics, or whether the presence of distinct diffusion scales may allow for emergence and stability of Turing patterns and their coexistence with jump-discontinuity patterns.

Systems coupling diffusive and non-diffusive components arise naturally in the mathematical modeling of spatially structured biological \cite{Klika12, Takagi21_2, He19}, chemical, and ecological processes \cite{Iuorio2023, marasco2014vegetation}. 
In many applications, only a subset of the variables diffuses in space, while other components remain spatially localized and evolve according to ordinary differential equations. 
This modeling framework captures interactions between spatially distributed extracellular species and localized variables. 
Typical examples include receptor-based signaling models \cite{Zhang2025}, pattern formation mechanisms in \emph{Hydra} \cite{Marciniak06}, and models of tumor growth and cancer signaling networks \cite{GALLAY2022103387}, where diffusive morphogens or signaling factors interact with intracellular regulatory dynamics confined to individual cells. In some cases reaction--diffusion--ODE models can be  rigorously derived from underlying microscopic descriptions using homogenization techniques \cite{marciniakptashnyk2008, marciniak2012stronghomog}.


Mathematically, such models consist of reaction--diffusion equations governing the evolution of spatially distributed species coupled with ordinary differential equations describing localized variables. 
The coexistence of diffusive and nondiffusive components introduces spatial scale separation that significantly influences the dynamics of the system. 
In contrast to classical reaction--diffusion systems, where pattern formation is typically analyzed through its {\it close-to-equilibrium} emergence mechanism, the interaction with nondiffusive variables may generate qualitatively different spatial structures. 
In particular, these systems may exhibit irregular or singular patterns, including spike solutions \cite{HMC14} or jump discontinuities \cite{Cygan2023, Akagi2024, Koethe20, Guo2025}, which are characteristic of {\it far-from-equilibrium} regimes.

A central question concerns \emph{de novo} pattern formation, namely the emergence of spatial structure from spatially homogeneous steady states. A classical mechanism underlying this phenomenon is diffusion-driven instability (DDI), also known as Turing instability \cite{turing1952chemical}.
DDI is a bifurcation from a stable spatially homogeneous steady state with respect to diffusion parameters, occurring when the effective diffusion ratio exceeds a critical threshold \cite{Nishiura1982}. It may lead to Turing patterns, which typically arise as small-amplitude, spatially heterogeneous perturbations of the homogeneous equilibrium near the onset of DDI.
 
In contrast, for systems coupling non-diffusive dynamics with a single reaction–diffusion equation, DDI may arise for arbitrarily small but non-zero diffusion. The  Reaction–diffusion–ODE systems thus exhibit a richer structure than classical reaction–diffusion models due to the presence of non-diffusive components, which introduce an additional spatial scale. Reaction-diffusion-ODE systems exhibit a richer structure than classical reaction-diffusion models due to the presence of non-diffusive components, which introduce an additional spatial scale.  

However, it has been shown that when only one component diffuses, all classical Turing patterns are unstable and the system cannot sustain regular stable heterogeneous steady states \cite{Cygan2022}. 
Hence, stable Turing patterns require models with at least two diffusive components.




Such models appear in applications and the diffusive components often exhibit what we call 
\textit{spatial scale separation}, meaning that the diffusive components evolve on 
substantially different spatial transport scales. This situation arises naturally in 
ecological models describing the interaction between vegetation biomass, soil water 
and toxic compounds accumulated in the soil \cite{marasco2014vegetation,Veerman2024_1}. In these 
models the transport mechanisms of the involved quantities differ significantly: soil 
water redistributes relatively fast through diffusion in the soil, plant biomass 
spreads much more slowly through vegetation growth and dispersal, while toxic 
compounds produced during plant decomposition act locally and are therefore often 
treated as nondiffusive variables. A model studied in \cite{Veerman2024_1} is given by
\begin{align}
D\,{ S_t} &= -S + B V + H V S , \\
 { V_t} &= \varepsilon^{2}\Delta V + U V^{2} - B V - H V S, \\
{ U_t} &= \Delta U + A(1-U) - U V^{2},
\end{align}
where $S$ denotes the toxic compounds, $V$ vegetation biomass and $U$ soil water. The 
model explicitly separates the spatial transport scales: water diffuses on the 
natural spatial scale, biomass diffuses much more slowly through the factor 
$\varepsilon^{2}$, and the autotoxic compounds are assumed to act locally and therefore 
do not diffuse. Consequently, 
the system naturally exhibits three spatial transport scales: fast diffusion, slow 
diffusion and the absence of spatial transport.

% {\color{red} Here maybe add something about Hydra. Extend}

Hair follicle patterning provides another natural example of such mechanisms, where receptor--ligand interactions involving the epidermal receptor Edar and diffusible morphogens such as BMPs, together with short-range mediators such as CTGF, form an activation--inhibition loop \cite{mou2006generation}. A corresponding minimal model, studied in \cite{klika2012influence}, takes the form
\begin{align}
\partial_t E &= F(E,B), \\
\partial_t C &= G(E,C) + D_C \Delta C, \\
\partial_t B &= H(E,C,B) + D_B \Delta B,
\end{align}
where $E$ denotes the nondiffusing receptor activity, $C$ the short-range diffusible component (CTGF), and $B$ the long-range diffusible inhibitor (BMP). The structural feature $D_C < D_B$ encodes the separation between local receptor dynamics, short-range transport, and long-range inhibition.

Another motivation for considering spatial scale separation comes from models with 
state-dependent diffusion. In the reduced 
Gatenby--Gawlinski model for tumor invasion studied \textit{e.g.} in~\cite{GALLAY2022103387} the spatial dynamics is governed by the system
% In this framework the variables $U$ and $V$ represent, respectively, the density of 
% healthy tissue and tumor cells. The spatial dynamics are governed by the system
\begin{align}
 U_t &= U\big( f(U) - dW \big), \\
 V_t &= \partial_x \big( f(U)\,\partial_x V \big) + rV f(V), \\
 W_t &= a\,\partial_x^{2} W + b(V - W),
\end{align}
where $U, V, W$ denote the density of healthy tissue, tumor cells and concentration of lactic acid, respectively. In this model the mobility coefficient depends on the local density of healthy tissue. More precisely, with $f(U)=1-U$, the effective mobility of the tumor population is reduced in regions where the healthy tissue remains close to its carrying capacity. Thus, although this mechanism is qualitatively different from diffusion with a small constant coefficient, it indicates that state-dependent diffusion may lead, in relevant regimes, to substantially weaker spatial spreading of one component than of another. From the modeling viewpoint, this provides a natural motivation for studying systems in which different components evolve on distinct spatial transport scales.

The first goal of this work is to understand how DDI manifests itself in reaction--diffusion--ODE systems with this three-scale structure. In particular, we derive simple and easily verifiable sufficient conditions for DDI in general systems with an arbitrary number of components. While the classical autocatalysis mechanism for the nondiffusive subsystem can still induce instability, it simultaneously destabilizes all nearby Turing patterns. This observation motivates the search for alternative mechanisms that can trigger DDI while still allowing for the existence of stable spatially heterogeneous steady states. Our analysis identifies such mechanisms and clarifies how interactions between nondiffusive, slow-diffusive, and fast-diffusive components can generate DDI compatible with Turing pattern formation.

In the multiple diffusive components setting, DDI may give rise to Turing patterns \cite{marasco2014vegetation,Lopez2025,Smith}. On the other hand, Turing patterns may also be unstable, and solutions may converge to far-from-equilibrium structures instead \cite{Marciniak06}. In particular, discontinuous stationary solutions exhibiting jump-type behavior may arise. Such patterns are analogous to those observed in reaction--diffusion--ODE systems with a single diffusive component, where far-from-equilibrium structures provide the dominant mechanism of spatial organization. In~\cite{Cygan2023}, a construction of discontinuous steady states based on a Banach fixed point argument was provided for general reaction--diffusion--ODE systems with a single diffusive component. This construction can be extended, with minor modifications, to systems involving multiple diffusive components. This forms the second goal of the present work. For the sake of completeness, we adapt this argument to our general framework and present it in a form suitable for the multi-scale setting considered in this paper.

The paper is organized as follows. Section \ref{sec:general_theory} presents sufficient conditions for DDI in general reaction--diffusion--ODE systems based on the decomposition into nondiffusive, slow-diffusive, and fast-diffusive components. Section \ref{sec:special_case} analyzes a three-component system in detail, where necessary conditions for DDI are derived and an example is provided in which instability is generated purely by the subsystem of diffusive species. Section \ref{sec:existence_jump} establishes the existence of far-from-equilibrium jump patterns. Finally, Section \ref{sec:Application} illustrates the theory using a receptor-based model and discusses implications for biological pattern formation, including morphogenesis in \emph{Hydra}.

\subsection{Mathematical setting}\label{sec:setting}
We consider a coupled reaction--diffusion--ODE system with $m = \mn + \ms + \mfa$ components of the form
\begin{align}\label{eq:full_system}
    \left\{
    \begin{aligned}
        \partial_t \un &= \fn(\un,\us,\uf),\\
        \partial_t \us &= \Ds \Delta \us + \fs(\un,\us,\uf),\\
        \partial_t \uf &= \Df \Delta \uf + \ffa(\un,\us,\uf)
    \end{aligned}
    \right.
\end{align}
posed in a bounded domain $\Omega\subset\RR^n$ with boundary $\partial\Omega\in C^2$.
Here, $\un$ denotes the $\mn$ nondiffusing components, $\us$ the $\ms$ slow-diffusing components, and $\uf$ the $\mfa$ fast-diffusing components. The diffusion matrices $\Ds$ and $\Df$ are diagonal with strictly positive entries, while the nondiffusing components have vanishing diffusion. Homogeneous Neumann boundary conditions are imposed on all diffusive species.

Let $F=(\fn,\fs,\ffa)$ denote the reaction terms, assumed to be at least $C^2$; weaker regularity suffices for our main results \cite{KMMü23}. The system can be linearized around a steady state $\NSS=(\Nun,\Nus,\Nuf)$ as
\begin{align}\label{eq:Model_lin_around_steady_state}
    \partial_t \xi &= \L \xi + \N(\xi), \\ \xi(0) &= \xi^0,
\end{align}
where $D:=\diag(0,\Ds,\Df)$ is the diffusion matrix, $J := \nabla_{(\un, \us, \uf)}F(\NSS)$ is the Jacobian at the steady state, $\xi$ denotes the perturbation, $\L$ the linearized operator given by
\[
\L \xi = D \Delta \xi + J \xi,
\]
and $\N(\xi)$ collects higher-order terms.

A constant steady state $\CSS=(\Cun,\Cus,\Cuf)$ undergoes diffusion-driven instability (DDI) if it is stable in the corresponding ODE system ($\Ds=\Df=0$) but loses stability when diffusion is introduced.
For reaction--diffusion--ODE models, nonlinear (in)stability of a bounded steady state follows directly from its linear (in)stability \cite[Theorems 3.3 and 3.7]{KMMü23}.
Linear stability reduces to the spectral analysis of the operator $\L=D\Delta+J$,
with spectral bound 
$$s(\L) = \sup_{\lambda\in\sigma(\L)} \Re \lambda.$$
For a bounded steady state $\NSS$, the spectrum of $\mathcal{L}$ on Banach space $L^p(\Omega)^{\mn+\ms+\mfa}$ or $L^\infty(\Omega)^{\mn}\times L^p(\Omega)^{\ms+\mfa}$ can be decomposed as
\begin{align}\label{eq:Spectrum_of_general_L}
    \sigma(\L) = \sigma\big(\nabla_{\un} \fn (\NSS)\big) \cup \sigma_p(\L),
\end{align}
under the assumption $\Ds, \Df >0$ \cite[Section 5.22]{Kowall_2}, \cite[Proposition 2.3]{KMMü23}.
In the case of a constant steady state $\CSS$ and $p=2$, the point spectrum is expressed in terms of the eigenvalues $(\lambda_j)_{j\ge0}$ of the Laplace operator $-\Delta$ with Neumann boundary conditions:
\begin{align}\label{eq:eigenvalues_at_constant_steady_state}
    \sigma_p(\L) = \bigcup_{j=0}^\infty \sigma(-\lambda_j D + J) = \bigcup_{j=0}^\infty \big\{\lambda\in\mathbb{C} \mid \det(\lambda I + \lambda_j D - J) = 0\big\}.
\end{align}
For later reference, we define the following submatrices of $J = \nabla_{(\un, \us, \uf)}F(\CSS)$:
\begin{equation}
\begin{aligned}
\label{eq:definition_submatrices}
    J_1 := \nabla_{\un} &\fn(\CSS)\in\mathbb{R}^{\mn\times\mn},\ J_2 := \nabla_{\us} \fs(\CSS)\in\mathbb{R}^{\ms\times\ms},\ J_3 := \nabla_{\uf} \ffa(\CSS)\in\mathbb{R}^{\mfa\times\mfa},\\
    J_{12} &:= \nabla_{(\un, \us)} (\fn,\fs)(\CSS) = \begin{pmatrix}
        \nabla_{\un} \fn(\CSS) & \nabla_{\us} \fn(\CSS) \\
        \nabla_{\un} \fs(\CSS) & \nabla_{\us} \fs(\CSS)
    \end{pmatrix} \in\mathbb{R}^{(\mn + \ms)\times(\mn + \ms)},\\
    J_{13} &:= \nabla_{(\un, \uf)} (\fn,\ffa)(\CSS) = \begin{pmatrix}
        \nabla_{\un} \fn(\CSS) & \nabla_{\uf} \fn(\CSS) \\
        \nabla_{\un} \ffa(\CSS) & \nabla_{\uf} \ffa(\CSS)
    \end{pmatrix} \in\mathbb{R}^{(\mn + \mfa)\times(\mn + \mfa)}, \\
    J_{23} &:= \nabla_{(\us, \uf)} (\fs,\ffa)(\CSS) = \begin{pmatrix}
        \nabla_{\us} \fs(\CSS) & \nabla_{\uf} \fs(\CSS) \\
        \nabla_{\us} \ffa(\CSS) & \nabla_{\uf} \ffa(\CSS)
    \end{pmatrix} \in\mathbb{R}^{(\ms + \mfa)\times(\ms + \mfa)}.
\end{aligned}
\end{equation}


\section{Conditions for diffusion-driven instability}\label{sec:general_theory}

In this section we derive sufficient conditions for DDI in the general setting \eqref{eq:full_system}. {Several mechanisms may lead to DDI, and in general the corresponding conditions on the structure of nonlinearities may be rather involved. The simplest situations arise when instability is already present in particular subsystems of the Jacobian matrix. Although this approach provides transparent sufficient criteria, it does not cover all possible scenarios; more subtle mechanisms may occur in higher-dimensional systems, see Remark~\ref{rem:DimN}}.

\subsection{General conditions for DDI}\label{sec:DDI_general_conditions}
We begin with a generalization of DDI condition based on instability of the non-diffusive subsystem previously established for systems with a single diffusive component  \cite{Cygan2022}. For systems with several diffusing components, an analogous condition follows from the spectral characterization \eqref{eq:Spectrum_of_general_L}.

\begin{proposition}[DDI inducued by the non-diffusive subsystem]\label{prop:DDI_autocatalysis}
Let $\CSS$ be a constant steady state satisfying $s(J) < 0$ and $s(J_1) > 0$. 
Then~$\CSS$ exhibits DDI for all diffusion matrices $\Ds, \Df > 0$. 
Furthermore, any continuous nonconstant stationary solution of system~\eqref{eq:full_system} intersecting $\CSS$ is unstable.
\end{proposition}

\begin{proof}
Since $s(J) < 0$, the steady state~$\CSS$ is stable with respect to spatially homogeneous perturbations. 
On the other hand, by~\eqref{eq:Spectrum_of_general_L} we have
$
s(\L) \geqslant s(J_1) > 0
$
for all $\Ds, \Df > 0$ and $\CSS$ is unstable with respect to heterogeneous perturbations.

Moreover if $\NSS$ intersects $\CSS$ at some point $x_0 \in \Omega$, then
\[
s(\tilde{\L}) \geqslant s(\nabla_{\un}\fn(\NSS)) \geqslant s(\nabla_{\un}\fn(\CSS)) = S(J_1) > 0,
\]
which shows instability of $\NSS$.
The above estimate of the spectrum of the multiplication operator is based on the spectral characterization
\[
\sigma\left( \nabla_{\un}\fn(\NSS)\right) = \bigcup_{x\in\Omega\setminus N} \sigma \left( \nabla_{\un}\fn(\NSS(x))\right),
\]
where $N\subset X$ is of zero-measure \cite{KMMü23} and the fact that $\NSS$ is continuous.
\end{proof}

\begin{remark}
The result does not require the presence of all three component types: nondiffusive, slow-diffusive, and fast-diffusive in the system \eqref{eq:full_system}. 
For Proposition~\ref{prop:DDI_autocatalysis}, it suffices to have at least one nondiffusive component and at least one diffusive component, i.e., $\mn \geq 1$ and $\ms + \mfa \geq 1$.
\end{remark}

The autocatalysis mechanism enforces DDI independently of the diffusion coefficients, but it additionally rules out the stability of Turing patterns (Proposition \ref{prop:DDI_autocatalysis}).
Hence, we are interested in a different mechanisms leading to DDI and possible Turing patterns.
Our goal is to identify alternative conditions that still induce DDI while remaining simple to verify in applications.
To this end, we adopt a three-scale framework distinguishing nondiffusive, slow-diffusive, and fast-diffusive components.
Within this setting, DDI may originate from intermediate subsystems rather than from the purely nondiffusive block.
We therefore consider the case where instability arises from the coupled block $J_{12}$ involving both the $\mn$ nondiffusive and the $\ms$ slow-diffusive variables rather than from the nondiffusive block $J_1$.
Writing the linearized operator $\L$ explicitly depending on the diffusion coefficients as $\L_{\Ds,\Df}$, we obtain:



\begin{theorem}\label{thm:transfer_DDI_positive_D1}
    If $s(\L_{0,\hat{D}})>0$ for some fast diffusion matrix $\hat{D}>0$, then there exists $\varepsilon>0$ such that $s(\L_{\Ds,\hat{D}})>0$ for all slow diffusion matrices $0\leqslant \Ds<\varepsilon$.
\end{theorem}

\begin{proof}
    Consider fixed diffusion coefficients $\hat{D} > 0$ such that $s(\L_{0,\hat{D}}) > 0$.
    First, we show that this implies the existence of an eigenvalue of $\L_{0,\hat{D}}$ with strictly positive real part.
    If such an eigenvalue exists directly, we proceed with the continuity argument.
    Suppose there exists no eigenvalue with positive real part.
    Hence, by~\eqref{eq:Spectrum_of_general_L}, there exists a spectral element $\lambda$ with $\Re \lambda > 0$, which belongs to $\sigma(J_{12})$.
    Then, by \cite[Section 2.1.2]{Klika12}, there exists a sequence of eigenvalues $\mu_j \in \sigma(J - \lambda_j D)$ such that $\mu_j \to \lambda$ as $j \to \infty$, with $\Re \mu_j > 0$ for sufficiently large $j$.
    Since $\mu_j$ is also an eigenvalue of $\L_{0,\hat{D}}$ for all $j\in\NN$ by \eqref{eq:eigenvalues_at_constant_steady_state}, this yields a contradiction.
    
    Next, we use the continuity of eigenvalues with respect to diffusion parameters.
    The point spectrum of $\L_{\Ds, \Df}$ is characterized by roots of the characteristic polynomial
    \begin{align}
        P(\lambda; \lambda_j, \Ds, \Df) = \det(\lambda I + \lambda_j D - J),
    \end{align}
    see \eqref{eq:eigenvalues_at_constant_steady_state}.
    Since polynomial roots depend continuously on their coefficients \cite{CuckerCorbalan}, eigenvalues vary continuously with $\Ds$ and $\Df$.
    Because there exists $k \in \mathbb{N}_0$ such that $P(\lambda; \lambda_k, 0, \hat{D})$ has a root with positive real part, 
    we can choose $\varepsilon > 0$ small enough so that for all $0 \leqslant \Ds < \varepsilon$, $P(\lambda; \lambda_j, \Ds, \hat{D})$ retains a root with positive real part.
\end{proof}

\begin{remark}[DDI induced by a mixed subsystem]
    Sufficient conditions for DDI follow directly from Theorem~\ref{thm:transfer_DDI_positive_D1}. In particular, if the Jacobian matrix $J$ is stable while the principal submatrix $J_{12}$ is unstable, then the constant steady state $\CSS$ exhibits DDI for any fixed $\Df>0$ and sufficiently small $\Ds \geqslant 0$. This shows that instability of the mixed subsystem $J_{12}$ alone is sufficient to induce DDI without destabilizing Turing patterns.

    Note that this result also applies to classical reaction--diffusion systems, where all components diffuse. The minimal assumptions are $\ms \ge 1$ and $\mfa \ge 1$.
\end{remark}

\begin{remark}\label{rem:s(L)>0}
    Combining the results of Theorem~\ref{thm:transfer_DDI_positive_D1} and equation~\eqref{eq:eigenvalues_at_constant_steady_state}, we conclude that
    \begin{align}
        s(\L) > 0 \quad \Longleftrightarrow \quad \exists\, j \in \mathbb{N}_0 : \; s(-\lambda_j D + J) > 0.
    \end{align}
\end{remark}

Finally, destabilization may also originate from other subsystems. While the previous results identify $J_1$ and $J_{12}$ as natural candidates for
inducing DDI, interactions within the block
$J_{23}$, corresponding to slow and fast diffusion, may likewise generate
instability for suitable diffusion rates (see
Example~\ref{ex:DDI_from_PDE_subsystem}). To better understand the different mechanisms, it is useful
to characterize situations in which DDI is
strictly excluded. According to the spectral decomposition in
equations \eqref{eq:Spectrum_of_general_L} and
\eqref{eq:eigenvalues_at_constant_steady_state}, diffusion cannot
destabilize the steady state if the following two conditions hold
simultaneously:
\begin{enumerate}[label=(\roman*),noitemsep,topsep=0pt]
    \item the nondiffusive submatrix $J_1$ is stable,
    \item the matrix $J-D$ remains stable for every diagonal matrix $D \geqslant 0$.
\end{enumerate}
Matrices satisfying the second condition are referred to as
\emph{strongly stable}. Using the concept of Volterra--Lyapunov
stability for matrices \cite{Cross}, one can formulate explicit criteria ensuring that
conditions (i) and (ii) hold, and therefore exclude DDI.

\begin{theorem}[Volterra--Lyapunov stability criterion]\label{thm:NoDDI}
    If $J$ is Volterra--Lyapunov stable, i.e.\ there exists a diagonal matrix $M>0$ such that $JM+MJ^T$ is a negative definite matrix, then the constant steady state $\CSS$ is stable for all diffusion coefficients~$\Ds, \Df$.
\end{theorem}

\begin{proof}
    The matrix $J$ satisfying $JM + MJ^T<0$ for some diagonal matrix $M>0$ is strongly stable and all submatrices of $J$ are semistable~\cite[Proposition~1, Theorem~1]{Cross}. By \eqref{eq:Spectrum_of_general_L} and \eqref{eq:eigenvalues_at_constant_steady_state}, we have $s(\L)\leqslant 0$ for all diffusion coefficients $\Ds, \Df$.
\end{proof}

Let us sumarize the results of this section. The constant steady state $\CSS$ exhibits DDI if $s(J)<0$ and one of the following conditions holds:
\begin{equation}
\begin{aligned}
\label{eq:TwoConditionsDDI}
(i) &\quad J_1>0, && D_{\us} \text{ arbitrary,} && D_{\uf} \text{ arbitrary,} \\
(ii) &\quad s(J_{12})>0, && D_{\us} \text{ sufficiently small,} && D_{\uf} \text{ fixed}, \\
(iii) &\quad s(J_{13})>0, && D_{\us} \text{ fixed,} && D_{\uf} \text{ sufficiently small.}
\end{aligned}
\end{equation}
For completeness, we also include the case $s(J_{13})>0$, which
corresponds to the mirrored situation obtained by interchanging the
roles of the slow- and fast-diffusive components $v$ and $w$, although
this case does not fall directly within our framework.
In the case $(i)$, all Turing patterns are unstable, whereas in case $(ii)$ and $(iii)$ it is possible to observe Turing patterns.

\begin{remark}\label{rem:DimN}
For systems of dimension \(m \leqslant 3\), strong stability is equivalent to
the stability of \(J\) and all its principal submatrices
\cite[Theorems~3--4]{Cross}. Hence DDI can occur only if at least one
principal submatrix of \(J\) is unstable. For \(m>3\), this equivalence
fails: stability of all principal submatrices does not imply strong
stability \cite[Example~3.1]{Satnoianu}, and DDI may arise even when
every lower-dimensional subsystem is stable.
\end{remark}

Remark \ref{rem:DimN} shows that, although $J_1$ and $J_{12}$ provide
transparent mechanisms for DDI,
destabilization in reaction--diffusion--ODE systems cannot in general
be reduced to low-dimensional subsystems. In particular, for systems
with $m>4$, a complete characterization of all mechanisms leading to
DDI becomes complicated. For $m=3$, however, the situation can be analyzed in greater detail.
In this case one can systematically investigate mechanisms inducing
DDI beyond the instability of $J_1$ and $J_{12}$. Accordingly, we classify all possible sources of
DDI in the minimal three-component setting
consisting of one nondiffusive, one slow-diffusive, and one
fast-diffusive variable.





\subsection{A minimal system coupling nondiffusive, slow-diffusive and fast-diffusive components}\label{sec:special_case} 



To fully characterize which subsystems can induce DDI, we employ the Routh--Hurwitz criterion. For three-component systems, the stability conditions can be expressed explicitly in this
framework. Following \cite{Sakamoto12}, let
\begin{align}
    P(\lambda) = \lambda^3 + p_1 \lambda^2 + p_2 \lambda + p_3
\end{align}
be the characteristic polynomial of $J$, with
\begin{align}\label{eq:p1_p2_p3}
p_1 &= -\tr J, &
p_2 &= \det J_{12} + \det J_{23} + \det J_{13}, &
p_3 &= -\det J.
\end{align}
The Routh--Hurwitz criterion establishes that $J$ is stable (i.e., all roots of $P(\lambda)$ lie in the left half-plane) if and only if
\begin{align*}
    p_1 > 0, \quad p_3 > 0, \quad p_1p_2 - p_3 > 0,
\end{align*}
which also ensures $p_2 > 0$. We assume stability of $J$, a necessary condition for DDI.

DDI requires $\CSS$ to be unstable under spatially heterogeneous perturbations, i.e., $s(\L) > 0$. 
By Remark \ref{rem:s(L)>0}, this is equivalent to the existence of $j\in \mathbb{N}_0$ such that $s(- \lambda_j D + J) > 0$, where $\lambda_j$ denotes $j$-th eigenvalue of Laplacian $-\Delta$. 

We consider the characteristic polynomial of $-\mu D + J$ for $\mu\geqslant 0$, i.e.,
\begin{align}\label{eq:CharactPoly_-muD+J}
    P(\lambda; \mu) = \det(\lambda I + \mu D - J) = \lambda^3 + p_1(\mu) \lambda^2 + p_2(\mu) \lambda + p_3(\mu),
\end{align}
with coefficients
\begin{align*}
    p_1(\mu) &= \tr(\mu D - J),\\
    p_2(\mu) &= \det(J_{12} - \mu D_{12}) + \det(J_{23} - \mu D_{12}) + \det(J_{13} - \mu D_{13}),\\
    p_3(\mu) &= \det(\mu D - J).
\end{align*}
The assumption that $J$ is stable implies $p_1(\mu) > 0$ for all $\mu \geqslant 0$ and hence by the Routh--Hurwitz criterion, instability of $\L$ occurs if either
\begin{align}\label{eq:Routh-Hurwitz_criterion}
    p_3(\mu) < 0 \quad\text{or}\quad p_1(\mu)p_2(\mu) - p_3(\mu) < 0
\end{align}
for some $\mu=\lambda_j$ eigenvalue of $-\Delta$. 
% In this case, $\CSS$ exhibits DDI.

\begin{remark}\label{rem:Turing-Hopf}
    Assuming that $J_1<0$, $\det J_{12} > 0$ and $\det J_{13} >0$, one verifies that
    \begin{align*}
        p_3(\mu) = p_3 + [(\det J_{13})\Ds + (\det J_{12})\Df] \mu - J_1 \Ds \Df \mu^2
        > 0 \quad \text{for all} \quad \mu \geqslant 0,
    \end{align*}
    which excludes the existence of zero eigenvalues of $\L$.
    This means if we have DDI assuming these inequalities, we have Turing instability caused by the destabilization of a pair of complex conjugated eigenvalues, often called Turing--Hopf bifurcation, wave instability or W-instability.

    We know that for example $\tr J_{12} > 0$ is sufficient for DDI and produces a Turing-Hopf bifurcation when additionally it holds $\det J_{12} > 0$.
\end{remark}

We now determine which subsystems of $J$ must be unstable for the
steady state $\CSS$ to exhibit DDI. In particular, neither $J_2>0$ nor $J_3>0$ can alone induce DDI.
\begin{theorem}\label{prop:Discussion_possibil_DDI_3_equs}
If $\CSS$ exhibits DDI, then $J$ is stable and at least one of the subsystems $J_1$, $J_{12}$, $J_{13}$, or $J_{23}$ is unstable.
\end{theorem}
\begin{proof}
By \cite[Theorem~1.2]{Sakamoto12}, if $J$ and all its one- and two-component
subsystems are stable, then the steady state remains stable for all
$\Ds,\Df>0$. Equivalently, when $J_1$, $J_{12}$, $J_{13}$, $J_{23}$, and $J$ are all stable,
one verifies that $p_3(\mu)>0$ and $p_1(\mu)p_2(\mu)-p_3(\mu)>0$ for all
$\mu\geqslant 0$. 
\end{proof}

From~\eqref{eq:TwoConditionsDDI} we know that instability of $J_1$,
$J_{12}$, or $J_{13}$ is sufficient to induce DDI, in accordance with the classical LALI (local activation,
long-range inhibition) principle.
The case in which instability originates solely from $J_{23}$ is therefore of independent interest.
First, this situation lies outside the classical Turing paradigm,
which assumes an unstable slow- or nondiffusing activator coupled
to a stable fast-diffusing inhibitor.
Second it will always cause a Turing--Hopf bifurcation, i.e.\ a Hopf bifurcation purely induced by the Turing mechanism, see Remark \ref{rem:Turing-Hopf}. 
Third, the instability of $J_{23}$ can induce DDI but is not a sufficient condition, see Example \ref{ex:DDI_from_PDE_subsystem}.
However, it is easy to check that $\det J_{23} < 0$ is crucial to obtain DDI as $\tr J_{23} > 0$ alone cannot change the sign of $p_1(\mu)p_2(\mu)-p_3(\mu)$.


In applications it is common to apply a quasi-steady-state
approximation (QSSA) to eliminate a nondiffusing component, under the
assumption that this reduction preserves the qualitative dynamics of
the system \cite{Haerting17,Zhang2024}. In the case of three equations, QSSA results in a
classical two-component reaction--diffusion system, in which DDI occurs precisely when the slow-diffusing component satisfies the autocatalysis condition.
However, the latter requires $\det J_{12} < 0$, what can be shown for $J_1 < 0$ using the implicit function theorem.  
Consequently, if DDI in
the full model is induced by instability of $J_{23}$ or by the
condition $\tr J_{12} > 0$, then it disappears after
applying the QSSA, see Example \ref{ex:DDI_from_PDE_subsystem}.
In such case, the Turing mechanism induces a Hopf bifurcation (see Remark~\ref{rem:Turing-Hopf}), a phenomenon that cannot occur in a system consisting of only two reaction--diffusion equations.
For a more detailed discussion under which conditions DDI is preserved after applying QSSA see \cite{Smith}.




% Often in application a quasi-steady-state approximation (QSSA) is used for a nondiffusing component to reduce the system, assuming that this does not change the overall behavior of the system. 
% In the case of three equations, the reduced system is a classical Turing system consisting of two diffusive equations. It is well-known that we have DDI in this system exactly when the slow-diffusing component satisfies the autocatalysis condition in the reduced system.
% However, one can prove that the reduced system has the autocatalysis condition only when $\det J_{12} < 0$ in the original full model using the implicit function theorem and assuming $J_1 < 0$.
% Hence, if we have DDI in the full system induced by $J_{23}$ or $\tr J_{12} >0$, then DDI is not preserved when QSSA is applied.
% Intuitively, this makes sense as in this case the Turing mechanism induces a Hopf bifurcation (Remark \ref{rem:Turing-Hopf}), a phenomenon which cannot occur in a system of two reaction--diffusion equations.


% The case $s(J_{23})>0$, corresponding to the subsystem formed by the slow- and fast-diffusive variables, is relevant in this context. After applying a quasi--steady--state approximation to the nondiffusive variable, the resulting reduced system involves only the diffusive components, and the onset of DDI is often inferred from the stability of this reduced diffusive subsystem. However, the qualitative stability properties of the full system need not be preserved under this reduction. It may happen that the ODE subsystem is stable and that the reduced PDE system obtained via the quasi--steady--state approximation is also stable, while the unreduced PDE--ODE system is unstable. In this situation the instability arises solely from the coupling between the nondiffusive and diffusive variables. This phenomenon is illustrated in the example below.


\begin{example}[DDI resulting form instability of $J_{23}$]\label{ex:DDI_from_PDE_subsystem}
    We choose
    \begin{align*}
        J = \begin{pmatrix}-1&9&1.5\\-9&-1&5\\-2&3.5&-1\end{pmatrix}.
    \end{align*}
    A direct computation yields
    \begin{align*}
        s(J) < 0, \quad J_1<0, \quad J_2<0, \quad J_3 < 0, \quad s(J_{12})<0, \quad s(J_{13})<0, \quad s(J_{23}) > 0
    \end{align*}
    and
    \begin{align*}
        p_1(\mu)p_2(\mu) - p_3(\mu) &= 3.75 + (-6.5 \Df + 71.5 \Ds) \mu\\
        &\quad+ (2\Ds^2 + 2\Df^2 + 6\Ds\Df) \mu^2 + (\Ds^2\Df + \Ds\Df^2) \mu^3.
    \end{align*}
    For instance, choosing $\Ds = 0.001$ and $\Df = 1$ gives
    \begin{align*}
        p_1(\mu)p_2(\mu) - p_3(\mu) < 0 \text{ for } \mu \in (0.767, 2.433).
    \end{align*}
    On $\Omega = (0,\pi)$ with $\lambda_1 = 1 \in (0.767,\, 2.433)$, we obtain $s(\L) > 0$, so the system exhibits DDI.

    
    Reducing the system to the two diffusive components, the steady state is stable and no DDI occurs.
    Indeed, consider the linear system
    \begin{align*}
        \partial_t \begin{pmatrix}\un\\\us\\\uf\end{pmatrix} = \begin{pmatrix}0\\\Ds \Delta \us\\\Df\Delta \uf\end{pmatrix} + \begin{pmatrix}-1&9&1.5\\-9&-1&5\\-2&3.5&-1\end{pmatrix} \begin{pmatrix}\un\\\us\\\uf\end{pmatrix}.
    \end{align*}
    Applying the quasi-steady-state approximation (QSSA) to eliminate $\un$ yields the reduced Jacobian
    \begin{align*}
        \bar{J} = \begin{pmatrix}-82&-8.5\\-14.5&-4\end{pmatrix}, \quad \tr\bar{J} = -86 < 0, \quad \det\bar{J} = 204.75 > 0.
    \end{align*}
Since $\bar{J}$ and its principal submatrices are stable, the reduced system is stable for all choices of diffusion coefficients $\Ds, \Df > 0$.
\end{example}





\subsection{Large fast diffusion and domain rescaling}

The conditions derived in Section \ref{sec:DDI_general_conditions} describe DDI in terms of the diffusion coefficients by fixing the diffusion coefficients of the fast-diffusing subsystem and taking the remaining diffusion coefficients sufficiently small.
An equivalent formulation can be obtained by fixing the slow-diffusion coefficients and changing the fast-diffusion coefficients. However, this would additionally require a rescaling of the spatial domain.

{From the perspective of applications, these viewpoints play
different roles. While some diffusion coefficients may vary, others are
typically fixed physical parameters and cannot be adjusted. In contrast,
the size of the spatial domain may change, for instance through growth
of biological tissue or through a change in the spatial scale of the
system under consideration. Consequently, certain instability mechanisms
may become observable only when the domain is sufficiently large.}


% The conditions derived above describe DDI in
% terms of the diffusion coefficients, fixing the diffusion coefficients of the fast-diffusive components and choosing sufficiently small diffusion coefficients for the slow-diffusing components.
% % we fix large diffusion and choose smaller diffusion sufficiently small.
% An equivalent formulation can be obtained when fixing the coefficients of the slow-diffusive components, choosing sufficiently large fast-diffusive and by expressing these conditions through the size of the spatial domain.
% % Indeed, fixing the smaller diffusion coefficients, choosing larger diffusion sufficiently large and then rescaling the domain by a factor $L$ changes the spectrum
% % of the Laplacian in the same way as rescaling the diffusion coefficients,
% % so that these two perspectives are mathematically equivalent.
% From the viewpoint of applications, however, these interpretations play
% different roles. Some diffusion coefficients are typically fixed physical
% parameters and cannot be varied, whereas the size of the spatial domain
% may change, for instance through growth of the biological tissue or
% variation of the spatial scale of the system under consideration. In
% such situations certain instability mechanisms may become observable
% only when the domain is sufficiently large. The following result
% reformulates the diffusion conditions obtained earlier in terms of
% domain rescaling and shows how sufficiently large fast diffusion can
% induce DDI on appropriately scaled domains.


% Rather than considering small slow diffusion rates $\Ds$, we may
% therefore equivalently examine the effect of large fast diffusion
% coefficients $\Df$. 

A simple rescaling argument yields the
following counterpart of Theorem~\ref{thm:transfer_DDI_positive_D1}.


\begin{corollary}
\label{cor:transfer_stab_incr_domain}
Let $\Ds > 0$ be fixed and suppose $s(\L_{0,\hat{D}}) > 0$ for some $\hat{D} > 0$.
Then, there exists $d_0 \geqslant 1$ and $L_0 \geqslant 1$ such that
\begin{align}
s(\L_{\Ds,\,d\hat{D}}) > 0
\quad\text{for all}\quad d \geqslant d_0,
\end{align}
provided $\L_{\Ds,\,d\hat{D}}$ is considered on the scaled domain $L\Omega$ with $L \geqslant L_0$.
\end{corollary}

\begin{proof}
The claim follows by rescaling the system, applying Theorem~\ref{thm:transfer_DDI_positive_D1}, and then transforming back to the original variables.  
Since Theorem~\ref{thm:transfer_DDI_positive_D1} is stated for a fixed domain, the rescaling may result in a larger spatial domain $L\Omega$ with $L \geqslant 1$.

Assume $s(\L_{0,\hat{D}}) > 0$ for some $\hat{D} > 0$ and fix $\Ds = \mathrm{diag}(d_1,\dots,d_{\ms})$.  
By Theorem~\ref{thm:transfer_DDI_positive_D1}, there exists $\varepsilon > 0$ such that  
\[
s(\L_{D,\,\hat{D}}) > 0 \quad \text{for all} \quad 0 \le D < \varepsilon .
\]
Choose $\tilde{d} \ge 1$ so that $d_i / \tilde{d} < \varepsilon$ for every $i=1,\dots,\ms$.  
Then $s(\L_{\,\frac{1}{d}\Ds,\,\hat{D}}) > 0$ for all $d \ge \tilde{d}$.

Rescaling space by the factor $\sqrt{d}$ transforms the domain to
\[
\tilde{\Omega} = \sqrt{d}\,\Omega =: L_d \Omega
\]
and gives $s(\L_{\Ds,\,d\,\hat{D}}) > 0$ for all $d \ge \tilde{d}$.  
Since $\tilde{d} \ge 1$, we have $L_d \ge 1$, with $L_{\tilde{d}} = 1$ if $\varepsilon > 0$ is large enough to ensure $d_i < \varepsilon$ for all $i$.  
\end{proof}

\begin{remark}
\label{rem:increasing_domain_nec}
    Increasing only $\Df$ while keeping $\Ds$ fixed is, in general, insufficient to induce DDI. 
    An enlargement of the spatial domain with fixed $\Ds$ is required, since this effectively decreases $\Ds$ relative to $\Df$ and thereby restores the instability mechanism.

    In case of instability of $J_{12}$, we have $s(\L_{0,\hat{D}}) > 0$ for all $\hat{D} > 0$ and can choose $\hat{D}$ arbitrary.
    If additionally $J$ is stable, we have DDI for $\Ds$ fixed, $\Df = d\hat{D}$, $d>0$ sufficiently large, on domain $L\Omega$, $L\ge1$ sufficiently large.
\end{remark}

\begin{remark}
In case of a system of three equations, one can explicitly determine the domain scaling required for DDI with unstable $J_{12}$ using the Routh--Hurwitz criterion.
We assume that
\[
s(J)<0,\quad J_1 < 0,\quad  \Ds \text{ fixed}, \quad \Df \text{ arbitrarily large}.
\]
Depending on whether $\tr J_{12} > 0$, $\det J_{12} < 0$, or both, one can estimate the leading-order terms of $p_3$ and $p_1p_2 - p_3$ in $\Df$ to determine the range of $\mu$ where instability of $-\mu D + J$ occurs, and thus choose $L$ such that the first positive Laplacian eigenvalue $\lambda_1$ on $L\Omega$ lies within this range.
From this, we obtain:
\begin{itemize}
    \item Case $\tr J_{12} > 0$:

    For $\varepsilon \in \big(0, \frac{\tr J_{12}}{2\Ds}\big)$, there exists $\Df > 0$ sufficiently large such that
    \begin{align}
        p_1(\mu)p_2(\mu) - p_3(\mu) < 0 \quad\text{for}\quad \mu \in \left(\varepsilon, \frac{\tr J_{12}}{\Ds} - \varepsilon\right).
    \end{align}
    Choosing $L \geqslant 1$ large enough so that $\lambda_1 < {\tr J_{12}}/{\Ds}$ 
    % for the first positive Laplacian eigenvalue $\lambda_1$ on $L\Omega$ 
    yields DDI.

    \item Case $\det J_{12} < 0$:

    For $\varepsilon \in \big(0, \frac{\det J_{12}}{2 J_1 \Ds}\big)$, there exists $\Df > 0$ sufficiently large such that
    \[
        p_3(\mu) < 0 \quad\text{for}\quad \mu \in \left(\varepsilon, \frac{\det J_{12}}{J_1 \Ds} - \varepsilon\right).
    \]
    Choosing $L \geqslant 1$ so that $\lambda_1 < {\det J_{12}}/({J_1 \Ds})$ yields DDI for large $\Df$.

    \item Case $\tr J_{12} > 0$ and $\det J_{12} < 0$:  
    DDI occurs for large $\Df$ if
    \[
        \lambda_1 \in \left(0, \max\left\{\frac{\tr J_{12}}{\Ds}, \frac{\det J_{12}}{J_1 \Ds}\right\}\right).
    \]
\end{itemize}
\end{remark}



\section{Existence of far-from-equilibrium patterns}
\label{sec:existence_jump}

As discussed in the introduction, reaction--diffusion--ODE systems exhibit a richer structure than classical reaction--diffusion models due to the presence of nondiffusive components.
In particular, such systems may give rise to irregular or singular patterns, including spike-type solutions or jump discontinuities, which we refer to as far-from-equilibrium patterns.
In this section, we focus on the latter case.

We establish the existence of nonconstant stationary solutions to the reaction--diffusion--ODE system~\eqref{eq:full_system}, consisting of a nondiffusive block $\Nun$ coupled to diffusive components $(\Nus, \Nuf)$. Notably, the stationary solutions may exhibit discontinuities in $\Nun$. The corresponding stationary problem is given by
\begin{nalign}
\label{eq:StationaryProblem}
\left\{
\begin{aligned}
    0 &= \fn(\Nun,\Nus, \Nuf),\\
    0 &= \Ds \Delta \Nus + \fs(\Nun,\Nus, \Nuf), \\
    0 &= \Df \Delta \Nuf + \ffa(\Nun,\Nus, \Nuf),
\end{aligned}
\right.
\end{nalign}
where $\Ds > 0$ and $\Df > 0$ denote the diffusion coefficients of the slow and fast components, respectively.
Our goal is to construct \emph{far-from-equilibrium} patterns, that is, stationary solutions that differ from a spatially homogeneous steady state on a set of positive measure. The construction is based on the approach developed in \cite{Cygan2023}, suitably adapted to the present multi-scale framework.

\subsection{Two-branch structure in the \texorpdfstring{$\Nun$}{u}-nullcline}
Let $\CSS = (\Cun,\Cus,\Cuf)\in\RR^m$ denote a constant 
solution of~\eqref{eq:StationaryProblem} with Jacobian $J$ with submatrices as in~\eqref{eq:definition_submatrices}.

\begin{assumption}
\label{ass:StatSol}
The linearization matrices satisfy
\begin{nalign}
    \det J_1 \neq 0 \quad \text{and} \quad \det \left(J - D\lambda_j \right)\neq 0 \quad \text{for all Neumann eigenvalues } \lambda_j \text{ of } -\Delta \text{ on } \Omega.
    \end{nalign}
\end{assumption}
This ensures that the implicit function theorem applies to $\fn(\cdot, \us, \uf) = 0$, yielding a $C^2$ map
\begin{align*}
    \varphi:\mathcal{V}\times\mathcal{W} \to \mathbb{R}^{\mn},\quad \varphi(\Cus, \Cuf) = \Cun,\quad \fn(\varphi(\us,\uf),\us,\uf) = 0,
\end{align*}
for $(\us, \uf)$ in a neighborhood $\mathcal{V}\times\mathcal{W}$ of $(\Cus, \Cuf)$.
The key ingredient for our construction is the existence of \textit{multiple distinct solution branches} for the equation $\fn(\un, \us, \uf) = 0$ with respect to the nondiffusive variable $\un$, all defined on the same domain $\mathcal V \times \mathcal W \subseteq \RR^{\md}$.
\begin{assumption}[Two distinct branches]
\label{ass:DiscontinuousStationaryTwoBranches}
There exist two distinct $C^2$--maps $\varphi,\psi:\mathcal{V}\times\mathcal{W}\to\mathbb{R}^{\mn}$ with
\begin{align*}
    \fn(\varphi(\us,\uf),\us,\uf) = \fn(\psi(\us,\uf),\us,\uf) = 0,
\end{align*}
$\varphi(\Cus,\Cuf) = \Cun$ and $\varphi\neq\psi$ on $\mathcal{V}\times\mathcal{W}$.
\end{assumption}
These maps form two branches of the $\Nun$-nullcline and allow constructing stationary solutions of system~\eqref{eq:full_system} that switch between those branches.

\subsection{Localised branch switching}
We study the effect of spatially localised perturbations by allowing the reaction terms to differ on complementary subdomains of~$\Omega$. The main mechanism is that a modification of the kinetics restricted to a subset of sufficiently small measure can be treated as a perturbation of the original system.

We first establish uniform $W^{2,p}$-bounds for the diffusive subsystem.

\begin{lemma}
	\label{DisLemma1} 
    Fix $p\in [2, \infty)$ and let $ D^d\in\RR^{\md\times \md}$ be a diagonal matrix with positive entries and $J^d \in \RR^{\md\times \md}$ satisfying 
    \begin{align}
    \det(J^d - D^d\lambda_j) \neq 0 \quad \text{for all} \quad j\in\NN_0.
    \end{align}
    Then, the operator 
    \begin{align}
        \L^d := D^d\Delta + J^d \quad \text{with the domain} \quad D\left(\mathcal L^d\right)=\left\{\nu\in W^{2,p}(\Omega)^{\md}:\ \partial_n \nu=0\text{ on }\partial\Omega\right\}
    \end{align} 
    is invertible and $\| (\L^d)^{-1} \xi\|_{W^{2,p}(\Omega)} \leqslant C \|\xi\|_{L^p(\Omega)}$ 
    for each $\xi \in L^p(\Omega)^{\md}$.
\end{lemma}

\begin{proof}
Let $\nu\in D(\mathcal L^d)$. By the standard elliptic estimate for the Neumann problem,
\begin{align}
\|\nu\|_{W^{2,p}(\Omega)^{m_d}}
\leqslant C\Big(\left\|\mathcal L^d \nu\right\|_{L^p(\Omega)^{m_d}}+\|\nu\|_{L^p(\Omega)^{m_d}}\Big).
\end{align}
Spectral characterization given by formula \eqref{eq:eigenvalues_at_constant_steady_state} with $\mn = 0$ yields $0\notin \sigma(\L^d)$ on $L^2(\Omega)^{\md}$.
Therefore, $\mathcal L^d:D(\mathcal L^d)\to L^p(\Omega)^{\md}$ is invertible and the inverse is bounded for each $p \in [1,\infty)$. The proof can be found \textit{e.g.} in \cite[Corollary 2.4]{KMMü23} for $p>n/2$ but it can be easily extended for $p\geqslant 2$.
Therefore, setting $\nu = (\L^{d})^{-1} \xi$ we obtain
\begin{align}
\left\|\left(\L^d\right)^{-1}\xi\right\|_{W^{2,p}(\Omega)^{m_d}}
\leqslant C\Big(\|\mathcal \xi\|_{L^p(\Omega)^{m_d}}+\left\|\left(\L^d\right)^{-1} \xi\right\|_{L^p(\Omega)^{m_d}}\Big) \leqslant C\|\mathcal \xi\|_{L^p(\Omega)^{m_d}}.
\end{align}
\end{proof}





Now we state a perturbation lemma allowing different reaction terms on two complementary subdomains.


\begin{lemma}
	\label{DisTheorem1}
	Let $p>n/2$, $p\geqslant 2$ and $q_1, q_2 \in C_b^2(\RR^{\md},\RR^{\md})$ {\rm (}\textit{i.e.}\ all derivatives up to order two are bounded{\rm )}. 
    Set $J_d:=\nabla_\nu q_1(\overline{\nu})$ and assume that for some $\overline{\nu} \in \RR^{\md}$ we have  
	\begin{nalign}
		q_1(\overline{\nu}) = 0 \quad \text{and} \quad \det\left(J^d - D^d \lambda_j\right)\neq 0 \quad \text{for each} \ j\in\NN_0,
	\end{nalign} 
    with a diagonal matrix $D^d>0$.\\
    Then, there exists $\varepsilon_0>0$ such that for for all $\varepsilon\in(0,\varepsilon_0)$ there exits  $\delta>0$ such that for any open subset $\Omega_1\subset \Omega$ with \mbox{$\Omega_2 = \Omega\setminus\overline{\Omega_1}$} and $|\Omega_2| \le \delta$, the problem
	\begin{nalign}
		\label{DisDefSingle}
		D^d\Delta \nu + q_1({\nu}) \mathbbm{1}_{\Omega_1} + q_2({\nu}) \mathbbm{1}_{\Omega_2} = 0 \quad \text{in}\ \Omega
	\end{nalign}
	admits a weak solution $ \nu\in W^{2,p}(\Omega)^{\md}$ with $\|\nu-\overline{\nu}\|_{W^{2,p}(\Omega)^{m_d}}\leqslant \varepsilon$. 
\end{lemma}


\begin{proof}
The argument is a straightforward vector-valued extension of the proof of the corresponding scalar result based on the Banach fixed point theorem \cite[Lemma 3.2]{Cygan2023}. We therefore indicate only the modifications required in the present setting.

Let $\L^d = D^d\Delta + J^d$.
After replacing \(\nu\) by \(\overline{\nu}+\omega\), equation \eqref{DisDefSingle} can be written in the equivalent form
\begin{align}
\omega = \mathcal T_1(\omega)+\mathcal T_2(\omega),
\end{align}
where
\begin{align}
\mathcal T_1(\omega):=\left(\mathcal L^d\right)^{-1}\Bigl(\bigl(J\omega-q_1(\overline{\nu}+\omega)\bigr)\mathbbm{1}_{\Omega_1}\Bigr),
\qquad
\mathcal T_2(\omega):=\left(\mathcal L^d\right)^{-1}\Bigl(\bigl(J\omega-q_2(\overline{\nu}+\omega)\bigr)\mathbbm{1}_{\Omega_2}\Bigr).
\end{align}
We consider the ball
\begin{align}
B_\varepsilon:=\{w\in W^{2,p}(\Omega)^{\md}:\|w\|_{W^{2,p}(\Omega)^{\md}}\le \varepsilon\}.
\end{align}
Since $q_1$ is bounded and $q_1(\overline{\nu}) = 0$, Taylor theorem yields
    \begin{align*}
        \left|J_d\omega(x) - q_1\big(\overline\nu+\omega\big)\right|
        &\leqslant \left| \frac{1}{2} \sum_{j,l=1}^k \partial_{\nu_l\nu_j} q_1(\overline\nu) \omega_j(x) \omega_l(x)\right|
        \leqslant C\left|\omega(x)\right|^2,
\end{align*}
for $\|w\|$ sufficiently small. Using the boundedness of $(\mathcal L^d)^{-1}$ and the embedding
\(W^{2,p}(\Omega)\hookrightarrow L^\infty(\Omega)\), valid since \(p>n/2\), we obtain
\begin{align}
\|T_1(\omega)\|_{W^{2,p}(\Omega)^{\md}}
\le C\|\omega\|_{L^\infty(\Omega)^{\md}}^2
\le C\varepsilon^2,
\end{align}
and similarly
\begin{align}
\|T_1(\omega)-T_1(\widetilde \omega)\|_{W^{2,p}(\Omega)^{\md}}
\le C\varepsilon \|\omega-\widetilde \omega\|_{W^{2,p}(\Omega)^{\md}}
\end{align}
for \(\omega,\widetilde \omega\in B_\varepsilon\).

For the term on \(\Omega_2\), since $q_2$ is also bounded, we only use the Lipschitz bound
\begin{align}
|J_d\omega-q_2(\overline{\nu}+\omega)|\leqslant C(1+|\omega|),
\end{align}
which implies
\begin{align}
\|T_2(\omega)\|_{W^{2,p}(\Omega)^{\md}}
\leqslant C |\Omega_2|^{1/p}(1+\|\omega\|_{L^\infty(\Omega)^{\md}})
\leqslant C |\Omega_2|^{1/p}(1+\varepsilon),
\end{align}
and
\begin{align}
\|T_2(\omega)-T_2(\widetilde \omega)\|_{W^{2,p}(\Omega)^{\md}}
\leqslant C |\Omega_2|^{1/p}\|\omega-\widetilde \omega\|_{W^{2,p}(\Omega)^{\md}}.
\end{align}

Hence, choosing \(\varepsilon>0\) sufficiently small and then \(\delta>0\) such that
\(|\Omega_2|\leqslant\delta\), the map \(\mathcal T:=\mathcal T_1+\mathcal T_2\) defines a contraction on \(B_\varepsilon\).
\end{proof}

% \begin{proof}
% 	By mapping $\nu\to \nu+\overline{\nu}$, $q_1(\nu)\to q_1(\nu+ \overline{\nu})$ and $q_2(\nu)\to q_2(\nu+ \overline{\nu})$, we may assume without loss of generality that $\overline{\nu} = 0$. We write equation \eqref{DisDefSingle} in the equivalent form
% 	\begin{nalign}
% 		\label{T0}
% 		\nu &= \big(D^d \Delta  + \nabla_\nu q_1(0)\big)^{-1} \Big(\big( \nabla_\nu q_1(0) \nu - q_1(\nu)  \big) \mathbbm{1}_{\Omega_1} \Big) \\ &+ \big(D^d \Delta  + \nabla_\nu q_1(0)\big)^{-1} \Big(\big(\nabla_\nu q_1(0) \nu - q_2(\nu) \big) \mathbbm{1}_{\Omega_2}\Big) \\
% 		&= \mathcal{T}_1(\nu) + \mathcal{T}_2(\nu).
% 	\end{nalign}
%         By Lemma~\ref{DisLemma1} the operator $\big(D^d \Delta  +  \nabla_\nu q_1({0})\big)^{-1}$ with $J = \nabla_\nu q_1({0})$ satisfies the estimate
% 	\begin{nalign}
% 		\label{OperatorLNormW22Estimation}
% 		\left\| \big(D^d \Delta  +  \nabla_\nu q_1({0})\big)^{-1} \zeta \right\|_{W^{2,p}} \leqslant C\|\zeta\|_\Lp{}
% 	\end{nalign}
% 	for each $p\geqslant2$ and all $\zeta\in \Lp{}^{\md}$. 
% 	Our goal is to show that the right-hand side of equation \eqref{T0} defines a contraction on the ball 
% 	\begin{nalign}
% 		B_\varepsilon(0) = \left\lbrace \nu \in {W^{2,p}(\Omega)}^{\md}: \|\nu\|_{W^{2,p}} \leqslant \varepsilon \right\rbrace
% 	\end{nalign}
% 	provided $|\Omega_2|$ and $\varepsilon >0$ are sufficiently small. 


    
% 	First, we deal with the operator $\mathcal{T}_1$. Since {$ q_1 \in C_b^2(\RR^{\md},\RR^{\md})$} and $q_1(0) = 0$ by the Taylor expansion we get
%     \begin{align*}
%         \left|\nabla_\nu q_1(0)\nu(x) - q_1\big(\nu(x)\big)\right|
%         %&=\left| q_1\big(\nu(x)\big) - q_1(0) - \nabla_\nu q_1(0) \nu(x)\right|\\
%         &\leqslant \left| \frac{1}{2} \sum_{j,l=1}^k \partial_{\nu_l\nu_j} q_1(0) \nu_j(x) \nu_l(x)\right|
%         \leqslant C\left|\nu(x)\right|^2,
%     \end{align*}
%     where $|\nu(x)|$ denotes the euclidean norm of a vector $\nu(x) \in \RR^k$.
% 	Hence, by estimate \eqref{OperatorLNormW22Estimation} and the Sobolev embedding $W^{2,p}(\Omega) \subseteq L^\infty(\Omega)$ for $p>N/2$, there exists a constant $C >0$ such that for each $\nu\in B_\varepsilon(0)$ we have
% 	\begin{nalign}
% 		\label{OperatorT1BTBEstimation}
% 		\| \mathcal{T}_1(\nu) \|_{W^{2,p}} \leqslant C \left\| |\nu(\cdot)|^2\right\|_{L^\infty(\Omega)} \leqslant C \|\nu\|_{L^\infty(\Omega)}^2 \leqslant C\varepsilon^2 .
% 	\end{nalign}
% 	Analogously, by the Taylor expansion,
% 	\begin{nalign}
% 		\Big| \nabla_\nu q_1(0) \nu(x) &- q_1\big(\nu(x)\big)-\Big( \nabla_\nu q_1(0) \omega(x) - q_1\big( \omega(x)\big)\Big) \Big| \\ 
% 		& \leqslant \left| q_1\big(\nu(x)\big) - q_1\big(\omega(x)\big) +{\color{red}-} \nabla_\nu q_1\big(\omega(x)\big)\big(\nu(x)-\omega(x)\big)\right|  
% 		\\ &+ \left|\big( \nabla_\nu q_1\big(\omega(x)\big) -  \nabla_\nu q_1(0) \big)\big(\nu(x)-\omega(x)\big)\right| \\
% 		& \leqslant \frac{1}{2}\|\nabla_\nu^2 q_1 \|_{L^\infty} \big|\nu(x)-\omega(x)\big|^2  + \|\nabla_\nu q_1\|_{L^\infty} \big|\omega(x)\big|\big|\nu(x)-\omega(x)\big|. 
% 	\end{nalign} 
% 	Hence, as in inequalities \eqref{OperatorT1BTBEstimation}, for $\nu,\omega\in B_\varepsilon(0)$, we obtain
% 	\begin{nalign}
% 		\label{OperatorT1Contraction}
% 		\|\mathcal{T}_1(\nu) - \mathcal{T}_1(\omega)\|_{W^{2,p}} \leqslant C\varepsilon  \|\nu-\omega\|_\W2p.
% 	\end{nalign}
% 	Here, we choose $\varepsilon_0>0$ such that for all $\varepsilon\in (0,\varepsilon_0)$ it holds
% 	\begin{nalign}
% 		C\varepsilon^2 \leqslant \varepsilon/2 \quad \text{and} \quad C\varepsilon < 1/2.
% 	\end{nalign}
% 	This gives is the wanted estimates in equations \eqref{OperatorT1BTBEstimation} and \eqref{OperatorT1Contraction}.

% 	Now we turn to the operator $\mathcal{T}_2$. We begin by the following inequality 
% 	\begin{nalign}
% 		\label{OperatorT2PointwiseEstimate}
% 		\left| \nabla_\nu q_1(0)\nu(x) - q_2\big(\nu(x)\big)\right| \leqslant \| \nabla_\nu q_1\|_{L^\infty} \big|\nu (x)\big| + \|q_2\|_{L^\infty}.
% 	\end{nalign}
% 	Thus, by estimate \eqref{OperatorT2PointwiseEstimate} and the embedding $W^{2,p}(\Omega) \subseteq \Li$ {for $p>N/2$,} for each $\nu\in B_\varepsilon(0)$ we get
% 	\begin{nalign}
% 		\label{eq:DiscontinuesSolutionsOperatorT2BTB}
% 		\|\mathcal{T}_2(\nu)\|_{W^{2,p}} &\leqslant C \big( \|\nabla_\nu q_1\|_{L^\infty} \left\|  \nu \cdot  \mathbbm{1}_{\Omega_2}\right\|_{L^p} + \|q_2\|_\LiR \|\mathbbm{1}_{\Omega_2}\|_{L^p} \big) \\
% 		& \leqslant C|\Omega_2|^{\frac{1}{p}}\big(   \|\nu\|_{W^{2,p}} + 1 \big).
% 	\end{nalign}
% 	Analogously, we have 
% 	\begin{nalign}
% 		\label{OperatorT2DifferencePointiwseEstimate}
% 		\Big|\nabla_\nu q_1(0)\nu(x) &- q_2\big(\nu(x)\big) - \big( \nabla_\nu q_1(0) \omega(x) - q_2\big(\omega(x)\big) \big) \big|\\  & \leqslant \Big(\|\nabla_\nu q_1(0)\|_{L^\infty}+\| \nabla_\nu q_2(0)\|_{L^\infty}\Big) \big|\nu(x)-\omega(x)\big|.
% 	\end{nalign}
% 	Thus, by inequality \eqref{OperatorT2DifferencePointiwseEstimate},  as in estimate \eqref{eq:DiscontinuesSolutionsOperatorT2BTB}, for $\nu, \omega\in B_\varepsilon(0)$, we obtain
% 	\begin{nalign}
% 		\label{OperatorT2Contraction}
% 		\|\mathcal{T}_2(\nu) - \mathcal{T}_2(\omega)\|_{W^{2,p}} & \leqslant  C \|(\nu-\omega)\mathbbm{1}_{\Omega_2}\|_\Lp{} \\
% 		& \leqslant C |\Omega_2|^{\frac{1}{p}} \|\nu- \omega\|_{W^{2,p}}.
% 	\end{nalign}
% 	Here, we choose $\delta >0$ such that for $\Omega_2 \subseteq \RR^N$ satisfying $|\Omega_2|<\delta$ we have 
% 	\begin{nalign}
% 		C\big( \varepsilon + 1 \big)|\Omega_2|^{\frac{1}{p}} \leqslant \varepsilon/2 \quad \text{and} \quad C |\Omega_2|^{\frac{1}{p}} < 1/2.
% 	\end{nalign}
% 	The estimates \eqref{OperatorT1BTBEstimation}, \eqref{OperatorT1Contraction}, \eqref{OperatorT2PointwiseEstimate} and \eqref{OperatorT2Contraction} yield that $\mathcal{T}_1 + \mathcal{T}_2$ maps $B_\varepsilon(0)$ into itself and is a contraction on this ball.
%     Hence, the Banach fixed point argument completes the proof of the existence of a solution to equation \eqref{T0} in $B_\varepsilon(0)$.
% \end{proof}


\subsection{Main existence theorem} Combining the geometric assumptions on the $\Nun$-nullcline with the regularity and perturbation lemmas, we now establish the existence of far-from-equilibrium stationary solutions.
The result is obtained by applying the localised branch-switching lemma to the $(\Nus,\Nuf)$-subsystem and reconstructing $\Nun$ from the chosen nullcline branch.
In contrast to the conditions for DDI discussed in the preceding sections, the diffusion coefficients do not play a role for the existence as long as $D_v>0$, $D_w>0$ and Assumption \ref{ass:StatSol} is satisfied.

\begin{theorem}
	\label{DisExBan}
    Consider problem \eqref{eq:StationaryProblem} with a constant solution $\CSS$ satisfying Assumption~\ref{ass:StatSol} and Assumption \ref{ass:DiscontinuousStationaryTwoBranches}.
    There exists $\delta>0$ such that for any open subset $\Omega_1\subseteq\Omega$ with $\Omega_2 = \Omega \setminus \overline{\Omega_1}$ satisfying $|\Omega_2|< \delta$, problem \eqref{eq:StationaryProblem} admits a weak solution $(\Nus,\Nuf)= \big(\Nus(x),\Nuf(x)\big)$ to
		\begin{nalign}
			\label{eq:DinsontinousSolutionsVEquation}
			\Ds \Delta \Nus + \fs(\Nun, \Nus, \Nuf) = 0 \qquad \text{and} \qquad \Df \Delta \Nuf + \ffa(\Nus, \Nuf, \Nun) = 0, 
		\end{nalign}
		where
		\begin{nalign}
			\label{Uform}
			\Nun(x) = \begin{cases}
				\varphi\big(\Nus(x), \Nuf(x)\big), \quad x\in \Omega_1, \\
				\psi\big(\Nus(x), \Nuf(x)\big), \quad x\in \Omega_2,
			\end{cases}
		\end{nalign} 
		satisfies $\fn\big(\Nun(x),\Nus(x), \Nuf(x)\big) = 0$ for almost all $x\in \overline{\Omega}$. 
\end{theorem}

\begin{proof}%[Proof of Theorem \ref{DisExBan}] 
        Let us choose $\varepsilon >0$ so small that $B_\varepsilon(\Cus, \Cuf) \subseteq \mathcal{V}\times \mathcal W$.
        \\
        First, we restrict $\varphi$ and $\psi$ to $B_\varepsilon(\Cus, \Cuf)$ and then we extend them in arbitrary way to $\tilde{\varphi}, \tilde{\psi} \in C^2_b(\RR^{\md}, \RR^{\mn})$ such that $\tilde{\varphi}= \varphi$ and $\tilde{\psi} = \psi$ on $B_\varepsilon(\Cus, \Cuf)$. Now, we define 
        \begin{align}\label{eq:def_q1_q2}
            q_1(\us,\uf) = \begin{pmatrix}    
            \fs\left( \tilde{\varphi}(\us,\uf), \us,\uf \right)    \\
            \ffa\left( \tilde{\varphi}(\us,\uf), \us,\uf \right)
            \end{pmatrix} \qquad \text{and} \qquad 
            q_2(\us,\uf) = \begin{pmatrix}    
            \fs\left( \tilde{\psi}(\us,\uf), \us,\uf \right)    \\
            \ffa\left( \tilde{\psi}(\us,\uf), \us,\uf \right)
            \end{pmatrix},
        \end{align}
        which satisfy assumptions of Lemma ~\ref{DisTheorem1}. Obviously $q_1, q_2$ are $C^2$-bounded 
        % {\color{gray}need $C_b^2$ but this not clear as $g$ doesn't have to be bounded in $v$. so have to circumvent this by setting $q\equiv 0$ outside of $B_{2\varepsilon}(\Cus, \Cuf)$?}
        together with $q_1(\Cus, \Cuf) = 0$.
        Using $\fn(\varphi(v,w),v,w) = 0$ on $B_\varepsilon(\Cus, \Cuf)$, one can calculate 
        \begin{align}\label{eq:deriv_tilde_phi}
            \nabla_{(v,w)} \tilde\varphi(\Cus, \Cuf) =
            \begin{pmatrix}-\nabla_{\un}\fn\left(\CSS\right)^{-1} \nabla_{\us}\fn\left(\CSS\right) \\
            -\nabla_{\un}\fn\left(\CSS\right)^{-1} \nabla_{\uf}\fn\left(\CSS\right)
            \end{pmatrix}.
        \end{align}
        The same holds for $\nabla_{(\us,\uf)}\tilde{\psi}(\Cus,\Cuf)$. For notational simplicity, we write $\nabla f:= \nabla f(\CSS)$ and use similar abbreviations for other derivatives.

    Then, using {Assumption \ref{ass:StatSol}}, \eqref{eq:def_q1_q2}, \eqref{eq:deriv_tilde_phi} and applying Schur's formula for determinants of block matrices, we obtain
        \begin{nalign}
		\det \big(\nabla_{(\us,\uf)} & q_1(\Cus, \Cuf) - D^d\lambda_j\big) =        
        \frac{1}{\det\left(\nabla_u f\right)}\det\left(J - D\lambda_j\right)\neq 0.
	\end{nalign}    
	Hence for each $p\in[2,\infty)$, by Lemma \ref{DisTheorem1}, we obtain a solution $(\Nus,\Nuf)\in W^{2,p}(\Omega)^{\md}$ to problem~\eqref{DisDefSingle} which satisfies $\| \Nus-\Cus\|_{W^{2,p}} + \| \Nuf-\Cuf\|_{W^{2,p}} \leqslant \varepsilon$ for arbitrary small $\varepsilon>0$ provided $|\Omega_2|>0$ is sufficiently small. By the Sobolev embedding $W^{2,p}(\Omega) \subseteq L^\infty(\Omega)$ with $p>N/2$ 
	and $\varepsilon >0$ small enough, we obtain $\tilde{\varphi} = \varphi$ and $\tilde{\psi} = \psi$ for all $x\in \Omega$ which completes the construction of a solution to problem \eqref{eq:StationaryProblem}.
\end{proof}


\begin{remark}[Proximity to $\CSS$]\label{rem:Jump_pattern_far_from_CSS}
    The stationary solution $(\Nun, \Nus,\Nuf)$ constructed in Theorem \ref{DisExBan} stays close to the points 
    $\big(\Cun, \Cus,\Cuf\big)$  and 
    $\big(\psi (\Cus,\Cuf), \Cus,\Cuf\big)$
    in the following sense.
    There exists $\varepsilon_0>0$ such that, for every $\varepsilon \in (0,\varepsilon_0)$, we find $\delta(\varepsilon)>0$ such that the solution constructed in Theorem \ref{DisExBan} satisfies
    \begin{gather}
        \label{DisVar}
        \|\Nus - \Cus \|_{L^\infty}  < \varepsilon, \qquad \|\Nuf - \Cuf \|_{L^\infty} < \varepsilon, \quad \text{and} \\ 
        \|\Nun -  \Cun \|_{L^\infty(\Omega_1)} + \|\Nun - \psi (\Cus, \Cuf) \|_{L^\infty(\Omega_2)}  < C\varepsilon.
    \end{gather}  
    If such solution is stable, we refer to it as \emph{far-from-equilibrium patterns} since the solution is "far away" from the equilibrium $(\Cun, \Cus, \Cuf)$ on $\Omega_2$ (depending on the distance of the two branches of $f=0$).
\end{remark}

\begin{remark}[Continuity and jumps]
\label{rem:Conti}
By Sobolev embedding, the diffusive components $\Nus$ and $\Nuf$ belong to $C(\overline{\Omega})$ and are therefore continuous.  
In contrast, the nondiffusive component $\Nun$, defined by~\eqref{Uform}, may exhibit jump discontinuities across $\partial\Omega_1$ whenever $\varphi(\nu,\omega) \neq \psi(\nu,\omega)$ in $\mathcal{V}\times\mathcal{W}$.  
Such discontinuities occur when the parameter $\varepsilon>0$ in~\eqref{DisVar} is chosen sufficiently small.
\end{remark}

\begin{remark}[Multiplicity and stability]
The subset $\Omega_2\subset\Omega$ can be chosen arbitrarily, provided its measure is sufficiently small, yielding infinitely many distinct stationary solutions.  

The stability of these patterns is a separate issue and may depend on the geometry of $\Omega_2$ as well as on the diffusion coefficients.
In practice, $\Omega_2$ does not need to be very small, as illustrated in Section~\ref{sec:Application}.
However, in general, the constructed solutions are not guaranteed to be stable.

In particular, suppose that the submatrix $J_1 = \nabla_{\un} f$, associated with the nondiffusive variables and evaluated at either $\CSS$ or at $(\psi(\Cus,\Cuf),\Cus,\Cuf)$, has an eigenvalue with positive real part.
Then, for sufficiently small $\delta>0$, the stationary solutions constructed in Theorem~\ref{DisExBan} are linearly unstable. This follows from Remark~\ref{rem:Jump_pattern_far_from_CSS} together with the spectral characterization in~\eqref{eq:Spectrum_of_general_L} and standard perturbation arguments for linear operators.
% Stability of these patterns is a separate issue, potentially influenced by the geometry of $\Omega_2$ and the values of the diffusion coefficients.  
% In practice, $\Omega_2$ does not need to be very small, as illustrated in Section~\ref{sec:Application}.
% If the submatrix $J_1$ of the Jacobian $J$ at $\CSS = \big(\Cun, \Cus,\Cuf\big)$ or at $\big(\psi (\Cus,\Cuf), \Cus,\Cuf\big)$ is unstable, then all steady states obtained from Theorem \ref{DisExBan} are unstable for $\delta>0$ sufficiently small.
% This follows from Remark \ref{rem:Jump_pattern_far_from_CSS} and the spectral characterization in \eqref{eq:Spectrum_of_general_L}.
\end{remark}

\section{Application to a receptor-based model}\label{sec:Application}

In this section, we illustrate the theoretical framework developed above by applying it to a receptor-based model consisting of one ODE and two PDEs. The model is inspired by~\cite{Marciniak06}, and describes the dynamics of receptors (\(u\)), ligands (\(v\)), and enzymes (\(w\)) on a one-dimensional spatial domain.

Let $\Omega = (0,L)\subset\mathbb{R}$, with $L>0$. The governing equations are
\be \label{eq:toysys}
\left\{
\begin{aligned}
    \del_t u &= -\mu_1 u + m_1 uv (1 + u v)^{-1}, & \textup{for } \ x \in \bar\Omega, \ t > 0,\\
    \del_t v &= D_v \Delta v -\mu_2 v + m_2 uv (1 + uv)^{-1} - vw, & \textup{for } \ x \in \Omega, \ t > 0,\\
    \del_t w &= D_w \Delta w -\mu_3 w + m_3 uv (1 + uv)^{-1}, & \textup{for } \ x \in \Omega, \ t > 0,
\end{aligned}
\right.
\ee
supplemented with homogeneous Neumann boundary conditions for $v$ and $w$, and nonnegative initial data.
Here $\mu_i \in (0,1]$ denotes decay rates, the $m_i \in [2,+\infty)$ are production rates, and $D_v, D_w > 0$ are the diffusion coefficients of ligands and enzymes respectively. The receptor population, $u$, is assumed diffusion-free.
The nonlinear production terms follow a Michaelis--Menten type saturation in the ligand--receptor binding.


Introducing $X = \left(u, v, w \right)^T$, $X_0 = \left(u_0, v_0, w_0 \right)^T$, and $D = \diag(0, D_v, D_w)$, the system is written in vector form
\begin{align}
\del_t X = D \Delta X + F \bigl( X \bigr), \quad \text{for} \ x \in \Omega \ \text{and} \ t > 0, 
\end{align}
with nonlinearities
\[F(X)  = \bpm -\mu_1 u + m_1 uv (1 + uv)^{-1} \\ -\mu_2 v + m_2 uv (1 + uv)^{-1} - vw \\ -\mu_3 w + m_3 uv (1 + uv)^{-1}\epm = \bpm f(X) \\ g(X) \\ h(X) \epm. \]
The existence, uniqueness, and boundedness of solutions are established in Appendix~\ref{app2}.

For the analysis of DDI, we require that the reduced kinetic system
\begin{align}
    \dot{ X}(t) = F \bigl( X(t) \bigr), \qquad X(0) = X_0,
\end{align}
possesses an asymptotically stable, non-negative steady state. 
To simplify the notation, we introduce the quantities
\begin{align}
\label{eq:EtaDef}
\eta_1 := \frac{m_1}{\mf}, \quad \eta_2 := \frac{m_2}{\ml}, \quad \eta_3 := \frac{m_3}{\me}, \quad \alpha := m_2 \eta_1 - \eta_3.
\end{align}

To formulate the upcoming results, we first collect the parameter constraints into the following assumption.


\begin{assumption}\label{assump:Parameters}
Let $\mu_i \in (0,1]$, $m_i\ge2$ for $i=1,2,3$, and $\eta_1, \eta_2, \eta_3, \alpha$ as defined in~\eqref{eq:EtaDef}. We assume that the following three conditions hold
\begin{enumerate}
    \item\label{item:first_assumption} $\zeta > 0$, where $\zeta := \alpha - 2\mu_2,$
    \item\label{item:sec_assumption} $\Theta > 0$, where $\Theta := \zeta^2 - 4 \mu_2 (\mu_2 + \eta_3),$
    \item\label{item:third_assumption} $\alpha + 2\eta_3 + \sqrt{\Theta} > 2\left(\frac{\mu_1}{\mu_3} - 1\right)(\mu_2 + \eta_3)\frac{\eta_1 \eta_3}{m_1 + m_2}.$
\end{enumerate}
\end{assumption}

Under Assumption~\ref{assump:Parameters}, the kinetic system admits three nonnegative steady states, as stated below.

\begin{theorem}\label{thm:bistabletoysys}
Under Assumption~\ref{assump:Parameters} the associated kinetic system admits three constant steady states:
\begin{align}
\CSS_{0} := (0, 0, 0)
\quad\text{and}\quad
\CSS_\pm := (\bar u_\pm, \bar v_\pm, \bar w_\pm),
\end{align}
with
\begin{align}
\label{eq:steadystates}
\bar v_\pm= \frac{\alpha + 2\eta_3 \pm \sqrt{\Theta}}{2 \eta_1 (\mu_2 + \eta_3)},\quad
\bar u_\pm = \eta_1 - \frac{1}{\bar v_\pm},\quad
\bar w_\pm = \eta_3 - \frac{\eta_3}{\eta_1 \bar v_\pm}.
\end{align}
Here, $\CSS_{0}$ is stable, $\CSS_-$ is positive and unstable, and $\CSS_+$ is positive and asymptotically stable.
\end{theorem}

The proof of this result is given in Appendix \ref{app1}.
Assume from now on that the conditions of Theorem~\ref{thm:bistabletoysys} hold. 

\subsection{Diffusion-driven instability}
We now investigate whether the stable positive steady state $\CSS_+$ can destabilize through DDI.
Since $\CSS_{0}$ is Volterra--Lyapunov stable (Theorem~\ref{thm:NoDDI}) and $\CSS_-$ is unconditionally unstable, only $\CSS_+$ is relevant for DDI.

Let $J := J(\CSS_+)$ be the Jacobian matrix of the nonlinearities $F$ evaluated at $\CSS_+$. A transformation using the relation $F(\CSS_+) = 0$, allows us to express $J$ as follows:
\begin{align}    
J 
= \dfrac{1}{(1 + \bar u_+ \bar v_+)^2} \bpm -m_1 \bar u_+ \bar v_+^2 & m_1 \bar u_+  & 0 \\ m_2 \bar v_+ & -m_2 \bar u_+ ^2 \bar v_+ & -\bar v_+ (1 + \bar u_+ \bar v_+)^2 \\ m_3 \bar v_+ & m_3 \bar u_+  & -\me (1 + \bar u_+ \bar v_+)^2 \epm.
\end{align}

The diagonal element $J_1<0$ rules out the autocatalysis condition of Proposition~\ref{prop:DDI_autocatalysis}.
Furthermore, the $2\times 2$ submatrices $J_{13}$ and $J_{23}$ are stable (negative trace, positive determinant), leaving $J_{12}$ as the only viable source of DDI (Theorem~\ref{prop:Discussion_possibil_DDI_3_equs}).

\begin{theorem}[Instability of $J_{12}$]
\label{thm:J12unstable}
Suppose Assumption~\ref{assump:Parameters} holds and, in addition,
\begin{enumerate}
    \item[$(4)$] $\eta_1 < 2\left(\frac{\eta_3}{m_2} + \frac{2}{\eta_2}\right)$.
\end{enumerate}
Then $J_{12}$ has an unstable eigenvalue.
Consequently, there exists either $D_v > 0$ sufficiently small, or a pair $(D_w, L)$ with $D_w > 0$ and $L \geqslant 1$ both sufficiently large, such that $\CSS_+$ undergoes DDI.
\end{theorem}

\begin{proof}
From $\tr (J_{12}) = - \bar u_+ \bar v_+ (m_1 \bar v_+ + m_2 \bar u_+)$, it is clear that the sub-matrix $J_{12}$ has a negative trace. Thus, the $2 \times 2$ matrix $J_{12}$ has an eigenvalue with positive real part if and only if $\det J_{12}<0$. This occurs when the parameters satisfy
\begin{align}
    \det(J_{12}) 
    = m_1 m_2 \frac{\bar u_+ \bar v_+}{\big(1+\bar u_+ \bar v_+\big)^4} 
      \left( (\bar u_+ \bar v_+)^2 - 1 \right) < 0.
\end{align}
% holds. 

Clearly, this inequality holds whenever $\bar u_+ \bar v_+ < 1$.  By definition of $\bar u_+$, the inequality is satisfied when $\bar v_+ < 2 / \eta_1$. Substituting the explicit expression for $\bar v_+$, we obtain
\be
\label{eq:ineq}\alpha + \sqrt{\alpha^2 - 4 \mu_2 (\alpha + \eta_3)} <  2 \bigl( \eta_3 + 2 \mu_2 \bigr).\ee
Using the definition of $\alpha$, inequality \eqref{eq:ineq} simplifies to the desired condition
\mbox{$\eta_1 < 2 ( \frac{\eta_3}{m_2} + \frac{2}{\eta_2})$}.

If this inequality is satisfied, then $s(J_{12}) > 0$.  
By Corollary~\ref{cor:transfer_stab_incr_domain}, there exists either $D_v > 0$ sufficiently small, or a pair $(D_w, L)$ with $D_w > 0$ and $L \geqslant 1$ both sufficiently large, such that the steady state $\CSS_+$ exhibits DDI.
\end{proof}

Theorem~\ref{thm:J12unstable} guarantees the existence of DDI under broad parameter regimes.
The threshold values of $D_v,D_w$, and $L$ in Theorem~\ref{thm:J12unstable} can be characterized explicitly using the Routh--Hurwitz criterion applied to the characteristic polynomial of $J - \lambda_j D$ (Theorem~\ref{lem:DDI_toysys_exact}).
See Section~\ref{sec:special_case} for the notation.
For this system, it is straightforward to verify that the Routh--Hurwitz elements $p_1(\lambda_j)$, $p_2(\lambda_j)$, and $(p_1p_2 - p_3)(\lambda_j)$ are positive for all $\lambda_j \geqslant 0$ and $D_v, D_w > 0$. 
Instability can therefore arise only if $D_v$ and $D_w$ are chosen such that, for some fixed eigenvalue $\lambda_j$ the Routh--Hurwitz element $p_3(\lambda_j) := - \det \big(J - \lambda_j D\big)$ is negative.

\begin{theorem}
\label{lem:DDI_toysys_exact}
    Let $\Omega = (0, L)$ and consider the constant steady state $\CSS_+$ of~\eqref{eq:toysys}. This steady state exhibits DDI if either of the following conditions is satisfied:
    \begin{enumerate}
        \item[(I)] $L >0$, the diffusion coefficient $D_w > 0$ is fixed and 
        \begin{align} 
        0 < D_v < \varepsilon:=\sup_{j\in\NN}\dfrac{\det(J) - \det(J_{12}) \, D_w \lambda_j}{\lambda_j \lp \det(J_{13}) - J_1 \, D_w \lambda_j \rp}, 
        \end{align}
        \item[(II)] $D_v > 0$ fixed, and
        $$ L > \pi j \sqrt{\dfrac{J_1 \, D_v }{\det(J_{12})}}, \qquad D_w >  \dfrac{\det(J) - \det(J_{13})\,D_v \lambda_j}{\lambda_j(\det(J_{12}) - J_1\,  D_v \lambda_j)} > 0 $$
        holds for at least one $j \in \mathbb N$.
    \end{enumerate} 
\end{theorem}

\begin{proof}
    To determine the precise parameter regimes for DDI, we analyze the sign of the Routh--Hurwitz element $p_3(\lambda_j) = - \det(J - \lambda_j D)$. Instability requires $p_3(\lambda_j)<0$, and we distinguish the cases of small $D_v$ and large $D_w$ accordingly.
    
    By solving $p_3(\lambda_j) < 0$ for $D_v$, that is, finding values of $D_v$ such that
    \begin{align}
    -\det(J) + \biggl( \det(J_{13})\,D_v + \det(J_{12}) \,D_w \biggr) \lambda_j + \biggl( -J_1\, D_v D_w \biggr) \lambda_j^2 < 0,
    \end{align}
    we obtain condition~\emph{(I)}.  
    The right-hand side of the resulting inequality is positive for sufficiently large $j \in \mathbb{N}$ and converges to $0$, since $J_{12}$ is unstable while $J_1$ and $J_{13}$ are always stable.

    The second condition is obtained by fixing $j \in \mathbb{N}$ and rewriting $p_3(\lambda_j) < 0$ as
    \begin{align}
    \label{eq:toy_critical_d2}
        \lambda_j \big( \det(J_{12}) - J_1\, D_v \lambda_j \big) D_w 
        < \det(J) - \det(J_{13})\,D_v \lambda_j.
    \end{align}
    Since the right-hand side of~\eqref{eq:toy_critical_d2} is negative, this inequality can only hold for large $D_w$ if
    \begin{align}
        \det(J_{12}) - J_1 D_v \lambda_j < 0.
    \end{align}
    The latter can always be achieved by increasing the domain length $L$ if necessary. 
    Solving~\eqref{eq:toy_critical_d2} for $D_w$ and the above inequality for $L$ yields the two expressions stated in condition~\emph{(II)}.
\end{proof}

\subsection{Numerical simulations}
We first verify that the parameter constraints introduced in Assumption~\ref{assump:Parameters} and Theorem~\ref{thm:J12unstable} define a non-empty region. These conditions guarantee the existence of the asymptotically stable steady state $\CSS_+$ and the occurrence of DDI through $s(J_{12}) > 0$.

Let $p := (m_1, m_2, m_3, \mu_1, \mu_2, \mu_3)$ denote the vector collecting parameters of system~\eqref{eq:toysys}, and $\mathcal P := [2, +\infty)^3 \times [0, 1]^3$ the admissible parameter domain. Within $\mathcal P$ we define
\begin{align*}
\mathcal R_1 &:= \biggl\{ p \in \Pi : \zeta > 0, \, \text{ where } \zeta := \alpha - 2\mu_2\biggr\},\\
\mathcal R_2 &:= \biggl\{ p \in \Pi : \Theta > 0, \text{ where } \Theta := \zeta^2 - 4\mu_2(\mu_2 + \eta_3) \biggr\},\\ 
\mathcal R_3 &:= \biggl\{ p \in \Pi : \alpha + 2\eta_3 + \sqrt{\Theta} > 2 \lp \dfrac{\mu_1}{\mu_3} - 1 \rp (\mu_2 + \eta_3)\dfrac{\eta_1 \eta_3}{m_1 + m_2} \biggr\},\\
\mathcal R_4 &:= \biggl\{ p \in \Pi : \eta_1 < 2 \left(\dfrac{\eta_3}{m_2} + \dfrac{2}{\eta_2}\right) \biggr\},
\end{align*}
A numerical exploration shows that 
\begin{align*}
    p_\star = (2.5, 9.68, 7.0, 0.95, 0.95, 0.6) \in \bigcap_{i=1}^4 \mathcal{R}_i,
\end{align*}
confirming that the feasible region is non-empty.
Since the inequalities defining $\mathcal R_1,\dots,\mathcal R_4$ are strict, an open neighborhood $\mathcal{V}_{p_\star} \subset \mathcal P$ exists in which all constraints remain satisfied (see Figure~\ref{fig:paramspace}).


\begin{figure}[ht]
    \centering
    \includegraphics[width=1.0\textwidth]{figures_theo/figure3.pdf}
    \caption{Cross sections of the six-dimensional parameter space, $\mathcal R$, in the $(m_i, \mu_j)$-planes, $i, j = 1, 2, 3$.
    In each panel, the regions $\mathcal R_1,\dots,\mathcal R_4$ are superimposed.
    The overlap, enclosed by a black curve, corresponds to parameters satisfying all constraints simultaneously. 
    The black dot marks the feasible point $p_\star$.}
    \label{fig:paramspace}
\end{figure}

The bifurcation structure of system~\eqref{eq:toysys} can be characterized using the Routh--Hurwitz polynomial $p_3$, whose roots determine stability boundaries in the diffusion parameter plane. 
The Turing-unstable region consists of all pairs of diffusion coefficients for which at least one spatial mode becomes unstable.
For each spatial eigenmode $\lambda_j$, define
$$\Gamma_j := \biggl\{ (D_v, D_w) \in (\RR_{>0})^2 \ : \ p_3(\lambda_j; D_v, D_w) < 0 \biggr\}, \qquad j\in\NN,$$ 
where the dependence of $p_3$ on $D_v$ and $D_w$ is made explicit.
Here, $\partial \Gamma_j$ denotes the associated stability curve $p_3(\lambda_j; D_v,D_w)=0$. Each $\Gamma_j$ consists of diffusion pairs destabilizing the $j$-th mode. If $(D_v,D_w) \in \Gamma_{i_1}\cap\dots\cap\Gamma_{i_k}$, the linearized operator at $\CSS_+$ has at least $k$ unstable eigenvalues. 
The full Turing-unstable set is therefore the union
\begin{align*}
    \Gamma := \bigcup_{j=1}^\infty \Gamma_j.
\end{align*}

Any diffusion pair $(D_v, D_w) \in \Gamma$ induces DDI of the steady state~$\CSS_+$. Figure~\ref{fig:bifurcation_diagram} illustrates the structure of $\Gamma$ and its constituent regions.
This description permits selective pattern control: choosing $(D_v,D_w)$ within a prescribed $\Gamma_j$ perturbs the $j$-th spatial mode, so that, with suitable initial conditions, patterns with a predictable number of peaks can be generated. Representative simulations on $\Omega=(0,1)$ are shown in Figure~\ref{fig:sim_patterns}.

\begin{figure}[ht]
    \centering    \includegraphics[width=1.0\textwidth]{figures_theo/bifurcation_diagram.pdf}
    \caption{ Bifurcation diagram for system \eqref{eq:toysys} in the $(D_v, D_w)$-plane. The shaded region $\Gamma$ is the Turing unstable set, given by the union of all mode-specific regions $\Gamma_j$. \newline
    \textbf{Parameters}: $\mf=1.00$, $\ml=1.00$, $\me=0.60$, $m_1=2.50$, $m_2=9.68$, $m_3=7.00$.}
    \label{fig:bifurcation_diagram}
\end{figure}

\begin{figure}[ht]
\includegraphics[width=0.95\textwidth]{figures_theo/figure2.pdf}
    \caption{Patterned solutions to system \eqref{eq:toysys}.
    Each row corresponds to one component of $X=(u,v,w)$; each column shows the outcome for a different diffusion pair $(D_v,D_w)$. 
    In all cases, only a single spatial mode is unstable. \newline
    \textbf{Parameters}: $\mf=1.00$, $\ml=1.00$, $\me=0.60$, $m_1=2.50$, $m_2=9.68$, $m_3=7.00$. \newline
    \textbf{Init. Cond.}: $u_0 = \bar u + \xi$, \ \ $v_0 = \bar v + \xi$, \ \ $w_0 = \bar w + \xi$, \qquad $\xi = \frac{x}{10} \sin(10\pi x)$.} 
    \label{fig:sim_patterns}
\end{figure}



\subsection{Turing pattern vs.\ far-from-equilibrium patterns}

We distinguish between two fundamentally different types of spatial patterns. First, there are \textit{Turing patterns}, which
are
defined as regular structures that bifurcate from a constant steady state through DDI.
Their spatial scale is dictated by the unstable eigenmodes of the linearized operator.
In contrast, we observe \textit{far-from-equilibrium patterns}, {which} do not rely on DDI: they may exist as isolated solution branches, exhibit discontinuities or spikes, and form a continuum of stationary states depending on the nonlinear reaction terms.
These patterns do not exist in classical reaction--diffusion systems, but are intrinsic to mixed ODE--PDE dynamics.
Importantly, distinguishing between these pattern types can be challenging, as both types of patterns can evolve from a perturbation of a constant steady state exhibiting DDI. 


The patterns observed in our simulations of system \eqref{eq:toysys} highlight the connection between classical Turing instability and far-from-equilibrium dynamics. To understand their nature, recall the nullcline for the nondiffusive component, $\{(u, v) \in \RR^2 : f (u, v) = 0 \}$, consists of two branches:
\begin{align*}
    u = 0 \quad \text{and}\quad u = \frac{m_1}{\mu_1} - \frac{1}{v},
\end{align*}
and the constant steady state $\bar X_+$ exhibiting DDI lies on the nontrivial branch. 
When perturbed by small random noise, the emerging patterns are initially governed by the unstable eigenmodes of the Laplacian: if only the $j$-th eigenmode is unstable, the spatial structure behaves similarly to $\cos(j\pi x/L)$ with the corresponding number of peaks, see Fig.~\ref{fig:sim_patterns}.

Surprisingly, nonlinear interactions allow stable patterns associated with higher modes that are \emph{linearly stable}. Intuitively, 
Figure~\ref{fig:different_perturbations} shows that, when $\lambda_4$ is the only unstable mode, large-amplitude perturbations aligned with $\cos(5\pi x)$ or $\cos(6\pi x)$ in the nondiffusive component still yield patterns reflecting those modes. 
This demonstrates that nonlinearity can stabilize patterns which otherwise would be beyond the reach of classical Turing instability,
% \textcolor{blue}{beyond the reach of linear theory}
provided the perturbation is strong enough to escape the basin of attraction of the unstable eigenmode.

\begin{figure}[!htbp]
\centering
\begin{subfigure}{0.9\linewidth}
    \centering
    \includegraphics[width=0.9\linewidth]{figures_theo/pattern_l4_random_noise.pdf}
    \caption{Small random perturbation around $\CSS$:
    \newline
    \textbf{Init. Cond.}: $u_0 = \bar u + \xi$, \ \ $v_0 = \bar v + \xi$, \ \ $w_0 = \bar w + \xi$, \qquad $\xi = 0.01 \cdot \mathcal U(-1, 1)$
    }
\end{subfigure} \\
\begin{subfigure}{0.9\linewidth}
    \centering
    \includegraphics[width=0.9\linewidth]{figures_theo/pattern_l4_mimic_l5.pdf}
    \caption{Perturbation describing eigenmode $\lambda_5$:
    \newline
    \textbf{Init. Cond.}: $u_0 = \bar u + 0.1 \cdot \bar u \cos(5 \pi x)$, \ \ $v_0 = \bar v + \xi$, \ \ $w_0 = \bar w + \xi$, \qquad $\xi = 0.01 \cdot \mathcal U(-1, 1)$
    }
\end{subfigure} \\
\begin{subfigure}{0.9\linewidth}
    \centering
    \includegraphics[width=0.9\linewidth]{figures_theo/pattern_l4_mimic_l6.pdf}
    \caption{Perturbation describing eigenmode $\lambda_6$:
    \newline
    \textbf{Init. Cond.}: $u_0 = \bar u + 0.1 \cdot \bar u \cos(6 \pi x)$, \ \ $v_0 = \bar v + \xi$, \ \ $w_0 = \bar w + \xi$, \qquad $\xi = 0.01 \cdot \mathcal U(-1, 1)$
    }
\end{subfigure}
\caption{
Patterns from simulations with $D_v = 0.006$ and $D_w = 0.017$, 
for which only the eigenmode $\lambda_4$ is unstable. 
Different initial conditions select patterns associated with different eigenmodes. 
\newline
\textbf{Parameters}: $\mf=1.00$, $\ml=1.00$, $\me=0.60$, $m_1=2.50$, $m_2=9.68$, $m_3=7.00$.
}
\label{fig:different_perturbations}
\end{figure}


A main observation that distinguishes patterns presented in Figure \ref{fig:sim_patterns} and Figure \ref{fig:different_perturbations} from classical Turing patterns is that rather than stabilizing near the constant steady state, solutions drift toward the trivial branch $u\equiv0$ of the nullcline, where they settle into stable far-from-equilibrium patterns. 
This branch-switching mechanism explains both the non-smooth structure of the $u$ state variable in Fig.~\ref{fig:sim_patterns} (despite their Turing-like appearance), and the discontinuous stationary states constructed in Fig.~\ref{fig:towerpattern1}. 
Thus, while unstable eigenmodes originating from DDI act as a selection mechanism, the observed patterns are ultimately far-from-equilibrium states determined by the structure of the nonlinear branch.
We recreate this property by varying initial conditions, see Figure \ref{fig:different_perturbations} for comparison.

Theorem~\ref{DisExBan} provides the theoretical foundation: for any sufficiently small subset $\Omega_2\subset\Omega$, one can construct stationary solutions in which $u$ switches between the two nullcline branches.
These patterns do not require DDI for their existence, only the presence of multiple solution branches satisfying Assumptions \ref{ass:StatSol}--\ref{ass:DiscontinuousStationaryTwoBranches}.
This yields uncountably many distinct far-from-equilibrium patterns, in stark contrast to the discrete set of Turing modes in classical reaction--diffusion systems. 
The continuum of possibilities arises because $\Omega_2$ may be chosen arbitrarily, producing a rich family of spatial organizations not captured by bifurcation from a steady state.
In fact, it is possible to construct patterns that are arbitrarily close in the $L^{2}$-norm. Consequently, none of these patterns can be \emph{asymptotically} stable in $L^{2}$, since for every such pattern one can find an arbitrarily close initial condition that is itself a stationary solution. In contrast, the same patterns are not close in the $L^{\infty}$-norm, which allows for the possibility of asymptotic stability in $L^{\infty}$.

Figure~\ref{fig:towerpattern1} illustrates this mechanism explicitly. 
Starting from a continuous pattern obtained near DDI, we force $u$ to jump to the trivial branch on small subsets of the domain. 
The resulting stationary profiles exhibit genuine jump discontinuities in $u$, yet remain stable under time evolution. 
This shows that discontinuous patterns are not artifacts of numerics but inherent to the system’s dynamics.

\begin{figure}[ht!]
    \centering
    \begin{subfigure}{0.9\linewidth}
        \centering
        \includegraphics[width=0.9\linewidth]{figures_theo/pattern_jump_tower_top.pdf}
        \caption{Initial profile forced to zero on $(0.49,0.51)$ produces two jump discontinuities at $x = 0.49$ and $x = 0.51$ (dashed black).
        }
    \end{subfigure}
    \begin{subfigure}{0.9\linewidth}
        \centering
        \includegraphics[width=0.9\linewidth]{figures_theo/pattern_jump_tower_bottom.pdf}
        \caption{Forcing $u=0$ on two subintervals $(0.35, 0.45)$ and $(0.55, 0.65)$ produces jumps at $x = 0.45$ and $x = 0.55$ (dashed black).
        }
    \end{subfigure}
    \caption{
    Simulation of System~\eqref{eq:toysys} with initial condition (dotted gray) taken from the pattern in Figure~\ref{fig:sim_patterns} ($\lambda_4$ unstable) and modified by introducing an artificial jump to zero as indicated in the subfigures. The resulting stationary profile (black) exhibits jump discontinuities.
    \newline
    \textbf{Parameters}: $\mf=1.00$, $\ml=1.00$, $\me=0.60$, $m_1=2.50$, $m_2=9.68$, $m_3=7.00$.
    }
    \label{fig:towerpattern1}
\end{figure}

In the case of stable Turing patterns, one would expect that choosing diffusion coefficients closer to the threshold for DDI would yield patterns of smaller amplitude, as the system is near to the bifurcation point. In such a regime, the patterns should arise purely on the nontrivial branch of the nullcline $f=0$ not using any far-from-equilibrium dynamics. However, this behavior is never observed in our simulations.

In conclusion, although the patterns in system~\eqref{eq:toysys} may resemble classical Turing patterns, their origin is fundamentally different: unstable eigenmodes associated with DDI act only as a catalyst, guiding the solution toward far-from-equilibrium states created by branch switching in the nondiffusive component.
This hybrid mechanism enriches the pattern formation landscape and highlights phenomena that cannot occur in standard reaction--diffusion systems.



\section*{Acknowledgment}
{This work is supported by the German Research Foundation (DFG) under Germany’s Excellence Strategy EXC 2181/1 - 390900948 (the Heidelberg STRUCTURES Excellence Cluster) and through the Collaborative Research Center 1324 (SFB1324, project B6). Additional funding was provided by the European Research Council (ERC) under the Synergy Grant (PEPS, No. 10107178)}


%%%%%%%%%% Insert bibliography here %%%%%%%%%%%%%%
% \newpage
% \vskip2pc
\bibliographystyle{abbrv}
\bibliography{biblio.bib}



\appendix


\section{Proof of Theorem \ref{thm:bistabletoysys}}
\label{app1}

We now prove Theorem \ref{thm:bistabletoysys}. The argument is divided into three steps. 
First, we show that the system admits three steady states when condition \ref{item:sec_assumption} holds.  
Second, we establish that the two nontrivial steady states are positive under condition \ref{item:first_assumption}.  
Finally, we analyze the Jacobian to determine that $\CSS_+$ is locally asymptotically stable, while $\CSS_-$ is unstable, under condition \ref{item:third_assumption}.

\begin{proof} \textbf{Step I: Analysis of steady states.} 
It is clear $\CSS_0 := 0$ is a constant steady state. At equilibrium, the relation $\frac{\del \bar u}{\del t} = 0$ implies $\bar u = \eta_1 - \frac{1}{\bar v},$ and substituting the expression of $\bar u$ into $\frac{\del \bar w}{\del t} = 0$ gives
\mbox{$\bar w = \eta_3 (1 - \frac{1}{\eta_1 \bar v}).$}
Plugging both relations into $\frac{\del \bar v}{\del t} = 0$ yields a quadratic equation that $\bar v$ must obey:
\begin{align}\label{eq:steadystatev}
    \eta_1(\mu_2 + \eta_3) \bar v^2 - (\alpha + 2\eta_3) \bar v + m_2 = 0.
\end{align}
The discriminant
\begin{equation}\label{eq:existssteadystate}
(\alpha + 2\eta_3)^2 - 4 m_2\eta_1 (\mu_2 + \eta_3) = (\alpha-2\mu_2)^2 - 4 \mu_2 (\mu_2 + \eta_3)
\end{equation}
is positive thanks to  assumption \ref{item:sec_assumption}.
Thus, equation \eqref{eq:steadystatev} has two distinct real solutions and overall, system \eqref{eq:toysys} admits three constant steady states: the trivial equilibrium $\CSS_0$,  and two nontrivial steady states $\CSS_\pm = (\eta_1 - \frac{1}{\bar v_\pm}, \bar v_\pm, \eta_3  - \frac{\eta_3}{\eta_1 \bar v_\pm}),$ where
\begin{align*}
\bar v_\pm := \dfrac{\alpha + 2\eta_3 \pm  \sqrt{(\alpha-2\mu_2)^2 - 4 \mu_2 (\mu_2 + \eta_3)}}{2 \eta_1 (\mu_2 + \eta_3)}.
\end{align*}

\vspace{1em}
\noindent
\textbf{Step II: Positivity of the nontrivial steady states.}
It is straightforward to verify $\bar v_\pm > 0$ for all admissible parameters.
From $\bar u_\pm = \eta_1 - 1/\bar v_\pm$ and $\bar w_\pm = \eta_3 (1 - 1/(\eta_1 \bar v_\pm))$, both $\bar u_\pm,$ and $\bar w_\pm$ are positive, if $\eta_1 \bar v_\pm > 1$.
Substituting the expression for $\bar v_\pm$ gives the inequality
\be\label{eq:condition_positive_+s}\mp \sqrt{(\alpha-2\mu_2)^2 - 4\mu_2(\mu_2 + \eta_3)} < \alpha - 2\mu_2.\ee
Recall that condition \ref{item:first_assumption} ensures $\alpha - 2\mu_2 > 0$. 
In the "$-$" case, inequality \eqref{eq:condition_positive_+s} holds automatically since the left-hand side is negative and the right-hand side is positive.
In the "$+$" case, both sides of equation \eqref{eq:condition_positive_+s} are positive, and the inequality \eqref{eq:condition_positive_+s} holds again. Therefore, both $\CSS_\pm$ are positive.

\vspace{1em}
\noindent
\textbf{Step III: Stability analysis.}
We investigate conditions under which the system is bistable. The Jacobian evaluated at an arbitrary equilibrium $\CSS = (\bar u,\bar v,\bar w)$ is
$$ \hspace{-2em} J(\CSS) = \dfrac{1}{(1+ \bar u \bar v)^2}
\bpm
-\mf (1+ \bar u \bar v)^2 + m_1 \bar v & m_1 \bar u & 0 \\
m_2 \bar v & -(\mf + \bar w) (1+ \bar u \bar v)^2 + m_2 \bar u & -\bar v (1+ \bar u \bar v)^2 \\
m_3 \bar v & m_3 \bar u & -\mu_3 (1+ \bar u \bar v)^2
\epm.$$
First, we observe $J(\CSS_0) = -\diag(\mf, \ml, \me),$ which is clearly stable.
However, more work is needed to conclude about the stability of $J$ around $\CSS_\pm$. Exploiting relations on the steady-states, we obtain
$$ J(\CSS_\pm) = \dfrac{1}{(1 + \bar u_\pm \bar v_\pm)^2} \bpm -m_1 \bar u_\pm \bar v_\pm^2 & m_1 \bar u_\pm  & 0 \\ m_2 \bar v_\pm & -m_2 \bar u_\pm ^2 \bar v_\pm & -\bar v_\pm (1 + \bar u_\pm \bar v_\pm)^2 \\ m_3 \bar v_\pm & m_3 \bar u_\pm  & -\me (1 + \bar u_\pm \bar v_\pm)^2 \epm.$$
Our next step is to use the Routh--Hurwitz criterion to show that $J$ is unstable around $\bar X_-$, and stable around $\bar X_+$. 

Computing $\det(J(\CSS_\pm))$ yields
\begin{align*}
\det(J(\CSS_\pm)) = \dfrac{1}{(1 + \bar u_\pm \bar v_\pm)^6} \biggl[ \bar v_\pm & (1 + \bar u_\pm \bar v_\pm)^2 \biggl(-m_1m_3 \bar u_\pm^2 \bar v_\pm^2 - m_1 m_3 \bar u_\pm \bar v_\pm \biggr) \\ & - \mu_3 (1 + \bar u_\pm \bar v_\pm)^2 \biggl( m_1 m_2 \bar u_\pm^3 \bar v_\pm^3 - m_1 m_2 \bar u_\pm \bar v_\pm \biggr) \biggr].
\end{align*}
Using the definition of $\bar u_\pm$, we establish $\bar u_\pm \bar v_\pm - 1 = \eta_1 \bar v_\pm - 2$. 
Thus,
\begin{align}
\label{eq:detJ}
\det(J(\CSS_\pm)) = -m_1 \mu_3 \dfrac{\bar u_\pm \bar v_\pm}{(1 + \bar u_\pm \bar v_\pm)^3} \biggl[\bigl(\alpha + 2 \eta_3 \bigr) \bar v_\pm  - 2 m_2\biggr].
\end{align}
For convenience, we introduce the notation
\begin{align*}
&\beta := \alpha + 2\eta_3 > 0, & &\gamma := 4\mu_2 (\alpha + \eta_3) > 0, & &\delta := 4 \eta_3 (\alpha + \eta_3) > 0.
\end{align*} 
Substituting the expression of $\bar v_\pm$ 
% from equation \eqref{eq:steadystates} 
into \eqref{eq:detJ} yields
$$-\det(J(\CSS_\pm)) = m_1 \me \dfrac{\bar u_- \bar v_-}{(1+\bar u_- \bar v_-)^3} \lp \dfrac{\beta}{2\eta_1 (\ml + \eta_3)} \lp \beta \pm \sqrt{\alpha^2 - \gamma} \rp - 2 m_2 \rp.$$
We stress that the square roots have positive argument, since $\alpha^2 - \gamma > 0$ by assumption \ref{item:sec_assumption}.
Factoring the $2\eta_1(\ml + \eta_3)$ term out, we obtain
$$- \det(J(\CSS_\pm)) = \dfrac{\mf \me}{2 (\ml + \eta_3) } \dfrac{\bar u_- \bar v_-}{(1+\bar u_- \bar v_-)^3} \lp \beta^2 - \delta - \gamma \pm \beta \sqrt{\alpha^2 - \gamma} \rp.$$
Using $\beta^2 - \delta = \alpha^2$, we see that $-\det(J(\CSS_\pm))$ has the same sign as
$$\mathcal A := \alpha^2 - \gamma \pm \beta \sqrt{\alpha^2 - \gamma} = \sqrt{\alpha^2 - \gamma} \lp \sqrt{\alpha^2 - \gamma} \pm \beta \rp.$$
We deal with the case "$-$" and "$+$" separately. In the "$-$" case, it holds
$$\mathcal A = \dfrac{\sqrt{\alpha^2 - \gamma}}{\sqrt{\alpha^2-\gamma} + \beta} \lp \alpha^2 - \beta^2 - \gamma \rp = - \dfrac{\sqrt{\alpha^2 - \gamma}}{\sqrt{\alpha^2-\gamma} + \beta} \lp \delta + \gamma \rp < 0,$$
so $-\det(J(\CSS_-)) < 0$ and the Routh--Hurwitz criterion enables us to conclude that $\CSS_-$ is unstable. On the other hand, in the "$+$" case, $\mathcal A > 0$ follows from \eqref{eq:existssteadystate}. Thus $-\det(J(\CSS_+)) > 0$.
Using the Routh--Hurwitz criterion once again, the steady state $\CSS_+$ is stable if
\be\label{cond:DDI}-\tr \bigl( J(\CSS_+) \bigr) \biggl(\sum_{1 \le i < j \le 3} \det \bigl( J(\CSS_+)_{ij} \bigr) \biggr) + \det \bigl( J(\CSS_+) \bigr) > 0.\ee
This last expression is rather involved, so we break it down into smaller parts. Using the positivity of all terms in $-\tr(J(\CSS_+))$, we obtain the (coarse) estimate $-\tr(J(\CSS_+)) > \mu_3$.
Next, by distributing $\frac{m_1}{(1 + \bar u_+ \bar v_+)}$ inside the parenthesis in equation \eqref{eq:detJ}, the coefficient $-\det(J(\CSS_+))$ is rewritten as
$$-\det(J(\CSS_+)) = \mu_3 \dfrac{u_+ v_+}{(1 + \bar u_+ \bar v_+)^2} \lp \dfrac{m_1 \eta_3 \bar v_+}{1 + \bar u_+ \bar v_+} - m_1 m_2 \dfrac{1- \bar u_+ \bar v_+}{1 + \bar u_+ \bar v_+}\rp.$$
Finally, a direct computation shows
$$\sum_{1 \le i < j \le 3} \det \bigl( J(\CSS_+)_{ij} \bigr) = \mu_3 \dfrac{u_+ v_+}{(1 + \bar u_+ \bar v_+)^2} \lp m_1 \bar v_+ + m_2 \bar u_+ + \eta_3 - \dfrac{m_1 m_2}{\mu_3} \dfrac{1- \bar u_+ \bar v_+}{1 + \bar u_+ \bar v_+} \rp.$$
Putting everything together, inequality \eqref{cond:DDI} holds if
$$m_1 \bar v_+ + m_2 \bar u_+ + \eta_3 > \dfrac{m_1 \eta_3}{(1 + \bar u_+ \bar v_+) \mu_3} \bar v_+.$$
Using the equality $1 + \bar u_+ \bar v_+ = \eta_1 \bar v_+$ with the estimate $\bar u_+ > \bar v_+$, we rewrite the last inequality as
\be\hspace{-3em} \bar v_+ > \lp \dfrac{\mu_1}{\mu_3} - 1 \rp\dfrac{\eta_3}{m_1 + m_2}.\ee
This inequality holds thanks to assumption \ref{item:third_assumption}.
Hence, $\CSS_+$ is locally asymptotically stable. Therefore, the kinetic system is bistable under our assumptions and the proof is completed.
\end{proof}



\section{Existence of a nonnegative, bounded, global solution to System (\ref{eq:toysys})}
\label{app2}

In this section, we show that the system \eqref{eq:toysys} has a global solution, which is nonnegative and bounded, using a bounded invariant region argument. Kowall et al.~\cite[Proposition C.2]{KMMü23} recently established that RD--ODE systems of the form \eqref{eq:toysys} admit a unique local-in-time mild solution.
We denote such a solution by $X = (u,v,w)$.

Our goal here is to extend this local solution to all times, by proving that solutions to \eqref{eq:toysys} remain \textit{a priori} bounded for all time.
This is achieved using the theory of bounded invariant regions (see ~\cite[Chap. 14, \S B]{Smoller1983}). We first recall the definition of an invariant region.

\begin{definition}[Invariant Region, \cite{Smoller1983}] A closed subset $\Sigma \subset \RR^3$ is called a (positively) invariant region for the initial value problem \eqref{eq:toysys} if any solution $X(x,t)$ with initial values in $\Sigma$ satisfies $X(x,t) \in \Sigma$ for all $x \in \Omega$ and for all $t \in [0, T_{\max}]$.
\end{definition}

It is well known that rectangular regions are particularly convenient for reaction--diffusion systems.
By \cite[Corollary 14.8, Theorem 14.11]{Smoller1983}, the hyper-rectangle 
$$\Sigma := \biggl\{ X = (u, v, w) \in \RR^3 \ : \ a_u \le u \le b_u, \quad a_v \le v \le b_v, \quad a_w \le w \le b_w \biggr\}$$
is invariant provided the nonlinear mapping $F$ points strictly into $\Sigma$ on $\del\Sigma$. 
If in addition the system \eqref{eq:toysys} has the $f-$stable property (see \cite[Definition 14.10]{Smoller1983}), it is sufficient that $F$ points non-strictly into $\Sigma$ on the boundary. 

The $f-$stability of system $\eqref{eq:toysys}$ follows by adapting the proof in \cite[Theorem 3]{Daub2013Nov}.
Thus, it is only necessary to verify that $F$ points inward on the boundary of $\Sigma$.

\begin{proof}
    Let $A_u, A_v, A_w > 0$ denote arbitrary constants, and define the rectangular region
    \begin{align*}
        \Sigma := \bigcap_{i=1}^3[0, A_i] \subset \RR^3.
    \end{align*}
    We introduce the six edge functions
    \begin{align*}
        G_1(u, v, w) &= -u, & G_2(u, v, w)&= -v, & G_3(u, v, w) &= -w, \\[0.7em]
        G_4(u, v, w) &= u - A_1, & G_5(u, v, w)&= v - A_2, & G_6(u, v, w) &= w - A_3,
    \end{align*}
    so that $\Sigma = \bigcap_{j=1}^6\{ G_j \le 0\}$.
    By definition, for any boundary point $(u_\star, v_\star, w_\star) \in \del\Sigma$, there exists a $j$ such that $G_{j}(u_\star,v_\star,w_\star) = 0$.
    The condition that the nonlinear mapping $F$ points into $\Sigma$ is 
    $$\grad G_{j}(u_\star, v_\star, w_\star) \cdot F(u_\star, v_\star, w_\star) \le 0.$$
    
    If $u_\star = 0$, $v_\star = 0$, or $w_\star = 0$, then $f(u_\star, v_\star, w_\star) = 0$, $g(u_\star, v_\star, w_\star) = 0$, or $h(u_\star, v_\star, w_\star) = 0$ respectively and the condition is satisfied. This already shows that solutions remain nonnegative. On the other hand, if $u_\star = A_1$, or $v_\star = A_2$ or $w_\star = A_3$, then $\grad G_{j} = 1$.
    Using the positivity of $u,v,w$ and the bound $\frac{uv}{1+uv} \le 1$, the following estimates follow
    \begin{align*}
    \hspace{-1em} f(A_1, v, w) \le m_1 - \mu_1 A_1, \quad 
    g(u, A_2, w) \le m_2 - \mu_2 A_2, \quad 
    h(u, v, A_3) \le m_3 - \mu_3 A_3.
    \end{align*}
    Therefore, by choosing $A_i \ge \eta_i$ for $i = 1, 2, 3$, it holds
    \begin{gather*}
    \hspace{-2em}\grad G_{4}(A_1, v, w) \cdot F(A_1, v, w) \le 0, \quad \grad G_{5}(u, A_2, w) \cdot F(u, A_2, w) \le 0,\\ \grad G_{6}(u, v, A_3) \cdot F(u, v, A_3) \le 0.
    \end{gather*}
    Hence $\Sigma$ is a bounded invariant region for system \eqref{eq:toysys}:
    any solution starting inside $\Sigma$ remains inside $\Sigma$ for all times. So the solution is also bounded, hence global. The proof is complete.
\end{proof}


% \vskip6pt

% \enlargethispage{20pt}


\end{document}
